\documentclass[11pt,letterpaper]{article}
\usepackage[T1]{fontenc}
\usepackage{lmodern,amsmath,amssymb,amsthm,mathtools}
\usepackage[margin=.82in,headheight=15pt]{geometry}
\usepackage{booktabs,longtable,array,microtype,enumitem,xcolor,fancyhdr}
\usepackage{listings,xurl,hyperref}
\definecolor{ink}{HTML}{15324D}
\definecolor{soft}{HTML}{EDF3F7}
\hypersetup{colorlinks=true,linkcolor=ink,urlcolor=ink,citecolor=ink,
 pdftitle={Smaller explicit Diophantine equations for Con(ZFC)},
 pdfauthor={AI-assisted research prepared for Vladimir Reshetnikov}}
\urlstyle{tt}
\pagestyle{fancy}\fancyhf{}
\fancyhead[L]{\small\color{ink}Diophantine consistency: graphs and shared digit bounds}
\fancyfoot[C]{\thepage}
\setlength{\parindent}{0pt}\setlength{\parskip}{5.5pt}
\setlength{\emergencystretch}{2em}
\setlist{itemsep=2pt,topsep=4pt}
\lstset{basicstyle=\ttfamily\scriptsize,breaklines=true,columns=fullflexible,
 keepspaces=true,showstringspaces=false,frame=single,rulecolor=\color{ink!25},
 backgroundcolor=\color{soft},aboveskip=8pt,belowskip=8pt}
\newtheorem{theorem}{Theorem}[section]
\newtheorem{lemma}[theorem]{Lemma}
\newtheorem{proposition}[theorem]{Proposition}
\newtheorem{corollary}[theorem]{Corollary}
\theoremstyle{definition}\newtheorem{definition}[theorem]{Definition}
\theoremstyle{remark}\newtheorem{remark}[theorem]{Remark}
\newcommand{\Z}{\mathbb Z}\newcommand{\N}{\mathbb N}
\newcommand{\Con}{\operatorname{Con}}\newcommand{\ZFC}{\mathsf{ZFC}}
\newcommand{\Refut}{\operatorname{Ref}}\newcommand{\Ht}{\operatorname{ht}}
\newcommand{\degp}{\operatorname{deg}}
\newcommand{\file}[1]{\path{#1}}
\newcommand{\Qeq}{\mathcal Q}
\newcommand{\Ueq}{\mathcal U}
\begin{document}
\thispagestyle{empty}
{\color{ink}\Large\bfseries Smaller explicit Diophantine equations for
\(\Con(\ZFC)\)}\par
\vspace{4pt}
{\large Graph energies, shared digit bounds, and Pell centers}\par
\vspace{9pt}
Prepared for Vladimir Reshetnikov and the ProveIt project\par
24 October 2026\par
\textit{AI-assisted mathematical research with exact reproducible artifacts.}

\begin{abstract}
We construct several explicit integer-coefficient polynomials whose
unsolvability, in their specified domains, is equivalent in ZFC to the
syntactic consistency of ZFC. The main conceptual improvement is a graph
compiler that checks the coefficient support actually used by a graph
energy, rather than requiring all possible pairs of factor monomials to
be distinct. Opposite occurrences are paired, compatible aliases are
identified, and a four-term quartic removes certain equality witnesses.
On exactly the source used by the previous unit-coefficient construction,
the resulting quartic falls from 2,215 variables and 30,468 monomials to
1,137 variables and 23,105 monomials. With the newer signed source, we obtain
unit-coefficient quartics in 951, 984, and 1,050 unrestricted integers.
Their constant term is one and every other nonzero coefficient is also
\(1\) or \(-1\).

A second construction pools binomial low-digit bounds by their common
denominator and combines this with a semantic Pell anchor. It gives a
quartic in 226 nonnegative integers, with 175 quadratic residuals and
10,314 monomials; a 227-variable companion has only 7,905 monomials.
A third, smaller improvement gives two ten-positive-variable polynomials
of degrees 323 and 318 with 87- and 88-operation evaluators.
All equations are literal finite objects in the archive. The new algebraic
and arithmetic arguments are proved here. Their consistency interpretation
uses explicitly identified, pinned ProveIt calibration theorems.
No global minimum, consistency proof, historical priority claim, or new
end-to-end proof-assistant verification is asserted.
\end{abstract}

\vfill
\begin{center}
\fcolorbox{ink!25}{soft}{\parbox{.91\textwidth}{\small
\textbf{Suggested starting points.}
For coefficient simplicity, inspect the 984-variable quartic with
every nonzero coefficient \(\pm1\).
For the direct proof-calculus route, inspect the 226-variable natural
quartic and its complete residual appendix.
For a short evaluator, inspect the ten-positive-variable formulas in
Section~\ref{sec:compact}.}}
\end{center}
\clearpage
\begingroup\small\setlength{\parskip}{0pt}
\tableofcontents
\endgroup
\newpage

\section{The results and their exact meaning}

Throughout, \(\N=\{0,1,2,\ldots\}\). Positive integers are written
\(\Z_{>0}\); they are never silently interchanged with \(\N\) or \(\Z\).
Let \(\Refut_{\ZFC}(p)\) be the pinned syntactic predicate saying that
\(p\) codes a finite refutation in ZFC, and put
\[
 \Con(\ZFC)\;:\Longleftrightarrow\;\neg\exists p\,\Refut_{\ZFC}(p).
\]
The target is a theorem \emph{about a fixed polynomial}, not a proof of
the truth of either side of that theorem.

\begin{theorem}[Delivered consistency equations]\label{thm:main}
For every equation in Table~\ref{tab:results}, the integer companion in
Section~\ref{sec:cap}, and the two compact equations in
Section~\ref{sec:compact}, let \(D\) be its stated coordinate domain and
\(n\) its number of unknowns. The new ordinary proofs in this report,
composed with the
identified inherited calibration, give
\[
 \ZFC\ \vdash\quad
 \bigl[\exists\mathbf x\in D^n\;P(\mathbf x)=0\bigr]
 \ \longleftrightarrow\ \exists p\,\Refut_{\ZFC}(p).
\]
Consequently ZFC proves
\[
 \boxed{\bigl[\text{there is no zero of }P\text{ in }D^n\bigr]
       \ \longleftrightarrow\ \Con(\ZFC).}
\]
For the graph equations, the new part is stronger: it is an unconditional
explicit bijection between their integer zero sets and the zero sets of
the pinned signed source systems.
\end{theorem}

\begin{table}[ht]
\centering\small
\begin{tabular}{@{}lrrrr@{}}
\toprule
Equation and domain & Variables & Degree & Monomials & Height\\
\midrule
Graph, fewest variables, \(\Z\) & 951 & 4 & 56,916 & 1\\
Graph, balanced choice, \(\Z\) & 984 & 4 & 30,076 & 1\\
Graph, sparse choice, \(\Z\) & 1,050 & 4 & 21,766 & 1\\
Graph, matched older source, \(\Z\) & 1,137 & 4 & 23,105 & 1\\
Direct calculus, \(\N\) & 226 & 4 & 10,314 & 87,839,583\\
Direct calculus, sparse, \(\N\) & 227 & 4 & 7,905 & 87,839,583\\
\bottomrule
\end{tabular}
\caption{Literal collected polynomials. Height is the largest absolute
value of a collected coefficient. All graph polynomials have constant one.}
\label{tab:results}
\end{table}

The additional compact equations have ten strictly positive unknowns,
degrees 323 and 318, and 87 and 88 live binary arithmetic operations,
respectively. Each has the same four fixed
large coefficients reconstructed from a complete finite recipe; these
are not extra variables. Its support is not expanded in this release.

\subsection{Which meaning of ``simplest'' improves?}

There is no established global simplicity minimum. The repository already
contains quartics with fewer unrestricted integer variables when large
coefficients are allowed, and higher-degree equations with fewer variables
still. We therefore compare within the relevant families and retain the
actual tradeoffs.

\begin{table}[ht]
\centering\small
\begin{tabular}{@{}lrrr@{}}
\toprule
Matched comparison & Variables & Monomials & Height\\
\midrule
Package 127, signed442 source & 2,215 & 30,468 & 1\\
New compiler, same signed442 source & 1,137 & 23,105 & 1\\
Package 129, natural233 & 233 & 50,762 & 88,674,690\\
New direct natural226 & 226 & 10,314 & 87,839,583\\
Package 129, natural234 & 234 & 9,019 & 87,741,134\\
New sparse natural227 & 227 & 7,905 & 87,839,583\\
\bottomrule
\end{tabular}
\caption{Comparisons to the inspected repository snapshot. All rows have
degree four. The final sparse comparison slightly increases coefficient
height; it is not a simultaneous improvement in every measure.}
\label{tab:comparison}
\end{table}

The same-source graph reduction saves 1,078 variables, about 48.7 percent,
and 7,363 monomials, about 24.2 percent. This isolates the new compiler
from the additional gain obtained by switching to the newer signed324
source. The 226-variable natural equation improves all displayed size
measures over natural233. Against natural234, the 227-variable alternative
has fewer variables and terms but a slightly larger coefficient height.

\begin{minipage}{\textwidth}
The comparison snapshot is ProveIt commit
\begin{center}\small
\texttt{1c5367452a0bd4dd7e9e4bceb1508b62dbd2ff61}.
\end{center}
The comparison includes the seven recently integrated packages 127--133.
Relevant references are listed at the end.
\end{minipage}

\section{A graph polynomial with exactly the source zero set}

This section is independent of Pell equations, Turing machines, and set
theory. It applies to a finite system of quadratic equations over \(\Z\).
The inherited signed324 source is a system
\[
 R_1(\mathbf z)=\cdots=R_{208}(\mathbf z)=0,
 \qquad \mathbf z\in\Z^{324},
\]
whose literal rows are retained as \file{graph/source324/equation.json}.
Its consistency calibration is discussed in Section~\ref{sec:calibration}.

\subsection{Normalization with a unique integral extension}

First replace nonunit integer coefficients by explicitly defined constant
and scale coordinates. For each positive coefficient \(c>1\), recursively
include the constants needed for
\[
 c=2\lfloor c/2\rfloor+(c\bmod2).
\]
The corresponding equation uses two occurrences of the half-constant,
each with coefficient \(-1\), rather than a coefficient \(-2\).
For \(c z_i z_j\), use a scale coordinate \(s_{c,i}=c z_i\), then
the product \(s_{c,i}z_j\). Linear and constant terms are handled by the
same fixed constants. The definitions are triangular and monic.

The result is a list of signed term occurrences
\[
 S_r=\sum_{j=1}^{s_r}\epsilon_{rj}m_{rj},
 \qquad \epsilon_{rj}\in\{-1,1\},\quad\deg m_{rj}\le2.
\]
Every original tuple has a unique extension to the new coordinates;
all added defining rows then vanish, and the other rows equal the
original residuals. For signed324, normalization gives 506 coordinates,
390 rows, and 1,748 signed occurrences. For signed442 it gives 650
coordinates, 498 rows, and 1,969 occurrences. Occurrences, rather than
only collected monomials, matter at this stage.

\subsection{Incidence energy and integral tree flows}

Let a finite directed graph have an integer-valued polynomial label
\(t_e\) on every edge. Define its divergence at vertex \(v\) by
\[
 d_v=\sum_{\operatorname{tail}(e)=v}t_e
       -\sum_{\operatorname{head}(e)=v}t_e,
 \qquad E_G=\frac12\sum_v d_v^2.
\]
The polynomial \(E_G\) has integer coefficients: an edge square occurs
at its two endpoints, while every cross term in a squared divergence
has an even coefficient. Parallel edges do not invalidate integrality.
They can invalidate a particular height-one bound, which is why the
collected coefficient check is retained.

\begin{lemma}[Tree completion]\label{lem:tree}
Partition the vertices into nonempty blocks. Inside every block the
free edges form a spanning tree and carry distinct unrestricted integer
variables. All other edge values are prescribed. There is an internal
flow with every divergence zero exactly when the prescribed divergence
sum in each block is zero. When it exists it is unique and integral.
\end{lemma}
\begin{proof}
Necessity follows by summing divergences within a block; internal edges
cancel. Root the internal tree. At each leaf, its one remaining tree-edge
value is determined, with coefficient \(1\) or \(-1\), by the other
incident values. Eliminate the leaf and proceed. The final root equation
is precisely the block sum. This proves existence, uniqueness and
integrality without division.
\end{proof}

Use one block for each normalized row and one common hub block. Place
prescribed signed monomial edges so that the row-block sums are exactly
the \(S_r\). The hub sum is then \(-\sum_rS_r\). Introduce free tree
flows inside the blocks. By Lemma~\ref{lem:tree},
\[
 \exists\text{ integral free flows } E_G=0
 \quad\Longleftrightarrow\quad S_r=0\text{ for every }r.
\]
The graph need not be a tree globally.

\subsection{Pairing opposite occurrences}

Suppose the same monomial occurs with sign \(+1\) in row \(r\) and
sign \(-1\) in row \(s\). Instead of two edges to the common hub,
use one edge from block \(r\) to block \(s\), with that monomial
as label. Each row sum is unchanged, and the two hub contributions
cancel. This removes a prescribed edge and its associated factor copy.

The delivered deterministic layouts match 137 opposite pairs on the
signed324 source and 141 on signed442. The initial layout avoids
parallel prescribed edges, giving an easy initial coefficient bound.
This is a construction choice, not a necessary hypothesis of the final
zero-set theorem. The independent verifier checks the final block sums
directly, so it need not trust the matching algorithm.

\subsection{Aliases and a unit path}

An \emph{alias} is a coordinate intended to equal a particular normalized
source coordinate. Initially, each prescribed monomial receives private
factor copies, with some shared second factors. Linear and constant terms
are made homogeneous of degree two using aliases of the number one.
Thus prescribed edge labels have degree two and free edge labels have
degree one. The copied source value of an alias is recorded explicitly.

For a class of \(k\ge3\) ordinary aliases define
\[
 C(a_1,\ldots,a_k)=\sum_{j=1}^k a_j^2
                 -\sum_{j=1}^k a_ja_{j+1},\qquad a_{k+1}=a_1.
\]
Since \(2C=\sum_j(a_j-a_{j+1})^2\), it is nonnegative over \(\mathbb R\)
and vanishes exactly when all its aliases agree. A singleton needs no
penalty. A two-element class can initially be extended by one dummy
alias to obtain a three-cycle.

For the unit aliases \(u_1,\ldots,u_k\), \(k\ge2\), use
\[
 \Ueq(\mathbf u)=1+\sum_{j=1}^k u_j^2-u_1-u_k
                    -\sum_{j=1}^{k-1}u_ju_{j+1}.
\]
Twice this polynomial is
\[
 (u_1-1)^2+\sum_{j=1}^{k-1}(u_j-u_{j+1})^2+(u_k-1)^2.
\]
It therefore vanishes precisely when all unit aliases equal one.
Ordinary classes, unit aliases, and free flows are disjoint.

\section{Why checking actual support permits much more compression}

\subsection{A global factor restriction is unnecessarily strong}

The earlier compiler guaranteed that every product of two prescribed
monomials uniquely identified its two occurrences. A sufficient condition
was the absence of a four-cycle in a bipartite factor graph. But the
incidence energy contains only products from edges sharing a vertex.
Products of nonadjacent edge labels never appear.

For a concrete example, put the edge labels
\(ab,ad,cd,cb\), in that order, on a directed four-cycle. Its energy is
\begin{align*}
 E={}&a^2b^2+a^2d^2+c^2d^2+b^2c^2\\
     &-a^2bd-acd^2-bc^2d-ab^2c.
\end{align*}
Every nonzero coefficient is \(1\) or \(-1\). Yet the factor graph has
a four-cycle, and \((ab)(cd)=(ad)(cb)\). Those two equal products come
from opposite edge pairs, so neither occurs in the energy. This is the
precise gap exploited by the new compiler.

\begin{theorem}[Certified alias identification]\label{thm:quotient}
Start with a graph construction with complete ordinary alias classes
and a unit class. Identify any coordinates only within the same class,
leaving flow coordinates distinct. Substitute those identifications into
the graph energy and rebuild valid equality penalties on the surviving
aliases: singleton classes need none, two-element classes use a dummy
three-cycle or Lemma~\ref{lem:quarticpair}, and the unit path retains at
least two aliases. If the complete collected polynomial has coefficients in
\(\{-1,0,1\}\), its coefficient height is one and its integer zero
set remains in explicit bijection with the original source zero set.
\end{theorem}
\begin{proof}
Every old zero already gives the same value to the aliases being
identified. Conversely, any tuple for the identified coordinates expands
by copying each coordinate into all of its old aliases. At a zero of the
new polynomial, nonnegativity makes every component vanish. The rebuilt
ordinary penalties enforce the source labels, the unit path enforces
one, and the graph divergences vanish. The row-block sums consequently
give every normalized equation. Tree completion supplies the unique
integral flows in the forward direction. All normalizing constants and
scales were also uniquely determined.

This is a zero-set argument; rebuilding a cycle need not give the same
polynomial as substituting into the old cycle. The energy substitution
and the rebuilt penalties must be checked separately. The height claim
is exactly the finite collected-coefficient condition in the statement.
It does not require global monomial-product uniqueness.
\end{proof}

The implemented search tries identifications of equal-labelled aliases
and updates every affected energy coefficient with exact integers. An
identification is accepted only when no collected coefficient leaves
\(\{-1,0,1\}\). There is no numerical sampling and no optimization
oracle in the definition of the final polynomial. A fixed profile and a
complete literal certificate are supplied; the verifier independently
checks the resulting object from the original source.

\subsection{A four-term equality removes a dummy coordinate}

\begin{lemma}[A unit-coefficient equality polynomial]\label{lem:quarticpair}
The polynomial
\[
 \boxed{\Qeq(a,b)=a^4+b^4-a^3b-ab^3
             =(a-b)^2(a^2+ab+b^2)}
\]
is nonnegative over \(\mathbb R\), vanishes exactly on \(a=b\), and
has degree four and four nonzero coefficients, all of height one.
\end{lemma}
\begin{proof}
The second factor equals \((a+b/2)^2+3b^2/4\). It is positive away
from \((0,0)\), which already lies on the diagonal. The displayed
factorization proves the assertions.
\end{proof}

For a two-alias class, replace its dummy three-cycle by \(\Qeq(a,b)\).
This saves one coordinate. Unlike a quadratic cycle, the new penalty
can overlap the graph's quartic terms, so we accept it only after checking
the full collected sum. The balanced equation saves 57 dummy coordinates
this way; the minimum-variable equation saves 58, the sparse equation 59,
and the matched older-source equation 67.

\begin{proposition}[A narrowly defined local optimum]
Among nonzero real-nonnegative polynomials in two variables with integer
coefficient height one and zero set exactly \(a=b\), degree four is
minimal. Four nonzero monomials are also minimal.
\end{proposition}
\begin{proof}
A globally nonnegative nonzero polynomial cannot have odd total degree:
its highest nonzero homogeneous part changes sign under simultaneous
negation. A polynomial of degree at most two vanishing on the entire
diagonal and nonnegative everywhere must be \(c(a-b)^2\), \(c>0\).
This also follows by writing the polynomial in \(a-b,b\); nonnegativity
forces the coefficient of the first-order transverse term to vanish.
Integrality gives \(c\ge1\), so its mixed coefficient has magnitude at
least two. Thus degree below four is impossible.

After substituting \(a=b=t\), the sum of the coefficients at each
total degree must be zero. With coefficient height one, every nonempty
degree group therefore has an even number of terms. With fewer than
four terms a nonzero candidate has exactly two opposite same-degree
monomials. On the positive quadrant their difference changes sign as
\(a/b\) passes through one. It cannot be globally nonnegative.
Lemma~\ref{lem:quarticpair} attains both bounds. This proposition is
about this local equality gadget, not global Diophantine simplicity.
\end{proof}

\subsection{The final equation and an exact lift identity}

Let \(\mathcal C\) be the surviving ordinary cycles and \(\mathcal T\)
the selected two-alias classes. Each delivered graph polynomial is
\begin{equation}\label{eq:graph}
 \boxed{P_G=\frac12\sum_v d_v^2
             +\sum_{C\in\mathcal C}C
             +\sum_{(a,b)\in\mathcal T}\Qeq(a,b)
             +\Ueq(\mathbf u)=0.}
\end{equation}
The fixed vertices, edges, edge labels and alias classes are completely
listed in its certificate. The polynomial JSON also lists every
collected coefficient directly. The factor \(1/2\) is not a hidden
rational coefficient in the collected polynomial: integrality and
height one are both checked coefficient by coefficient.

There is an explicit polynomial lift \(L\) from the original source
coordinates, obtained by computing its normalizing constants and scales,
copying the alias values, setting the units to one and filling tree
flows by leaf elimination. For arbitrary source tuples, not just zeros,
\begin{equation}\label{eq:lift}
 P_G(L(\mathbf z))=
 \frac12\left(\sum_r R_r(\mathbf z)^2+
                    \left(\sum_rR_r(\mathbf z)\right)^2\right).
\end{equation}
Indeed, only the block roots can retain divergence; they carry the
individual residuals and their negative sum. All equality penalties
vanish. In particular the right side is \emph{not generally just}
\(\sum R_r^2\). The independent checker proves the root-divergence
identities symbolically and checks the complete coefficient expansion.

\section{Pooling low binomial digits in the direct calculus construction}
\label{sec:direct}

The graph theorem is a general compiler improvement. The next construction
changes the actual arithmetic certificate for a finite C222 derivation.
It starts from package 129's fully specified 344-coordinate normalized
natural graph, whose semantic interface and proofs are retained in the
direct component. No K215 simulation and no Jones universal polynomial
is instantiated on this route.

\subsection{The inherited arithmetic interface}

The source has a horizon \(t\ge16\), a proposed factorial \(F\), and
parameters \(x,y\), with
\[
 R=2+t+x+y+F,\qquad K=2R,\qquad N=B^K,\qquad q=N^K.
\]
It also represents
\[
 V_0=(F^2+1)^t,\quad
 D_t=q^t,\quad D_{t+1}=qD_t,\quad
 D_y=q^{y+1},\quad D_x=q^{x+1}.
\]
These variable powers are \emph{semantic descriptions} of retained
quadratic Pell circuits. They are not operations of the final polynomial.
Their intended meaning is established in the inherited soundness order,
before the factorial/binomial interpretation is used.

There are nine independent positive CRT representatives
\(a,b,c,d,S,S',u,v,C\), and a derived nonnegative constructor label
\(J\). The Cantor constructor is
\[
 I(a,b)=2(a+b)(a+b+1)+4b.
\]
The actual label polynomial contains \(I(a,b)\) and other nonnegative
terms, so \(J\ge2a^2\). Eleven numerator requests have the following
four common low-digit denominators:
\begin{center}\small
\begin{tabular}{@{}ll@{}}\toprule
Numerator requests & Denominator\\\midrule
Factorial binomial and \(C\) & \(D_t\)\\
\(a,b,c,d,J\) & \(D_{t+1}\)\\
\(S,S'\) & \(D_y\)\\
\(u,v\) & \(D_x\)\\\bottomrule
\end{tabular}
\end{center}
Originally each remainder \(r_i\) had a private nonnegative slack
\(s_i\), with \(D_i=r_i+s_i+1\). We pool these eleven bounds into four.
The separate factorial digit-output condition \(q=B+o+1\) is retained.

\begin{lemma}[Shared low-digit budget]\label{lem:digits}
Let \(q\ge2\), \(k\ge1\), and \(n_1,\ldots,n_m\in\N\). Suppose
\[
 \sum_{i=1}^m\binom{n_i}{j}\le q-1\qquad(0\le j<k).
\]
For the canonical low residues
\(r_i=\sum_{j=0}^{k-1}\binom{n_i}{j}q^j\), one has
\(\sum_i r_i\le q^k-1\). Hence a single nonnegative slack satisfies
\[
 q^k=\sum_i r_i+s+1.
\]
Conversely this one row, with arbitrary natural remainders, restores
all old bounds by \(s_i=s+\sum_{h\ne i}r_h\).
\end{lemma}
\begin{proof}
Sum the coefficient bounds against \(q^j\) and use
\((q-1)\sum_{j<k}q^j=q^k-1\). The restoration identity is literal
and all its coefficients are nonnegative.
\end{proof}

Pooling strengthens the arbitrary-witness constraints. Separate bounds
\(r_i<q^k\) alone would not justify it. Completeness therefore needs a
new uniform digit-margin estimate, proved next.

\subsection{The actual radix has ample unused low-digit capacity}

This subsection is only a constructive \emph{completeness} argument.
Start from a finite C222 refutation and choose the inherited true
witnesses, so \(F=t!\) and the displayed powers have their genuine
values. Soundness will not assume these identities prematurely.

The pinned CRT construction chooses the representatives at most \(N^2\)
and then replaces only \(a\) by \(a+N^2\). Thus every representative
is at most \(2N^2\), and their sum is at most \(10N^2\). Since
\(B\ge2\),
\[
 N\ge4^R,\qquad \frac q{N^2}=N^{2R-2},\qquad N^2\ge6.
\]
Each lower index \(k\in\{t,t+1,y+1,x+1\}\) is at most \(R-1\).
For an ordinary representative and \(0\le j<k\),
\[
 \binom nj\le\max(1,n^j)
 \le(10N^2)^{R-2}<N^{2R-2}=q/N^2,
\]
because \(10^{R-2}<N^2\), already implied by \(N^2\ge16^R\).

For the factorial upper index \(n_0=F^2+t<R^2+R\), use
\(R^2+R\le2^{2R}\), \(t\le R\), and \(R\ge2\):
\[
 \binom{n_0}{j}\le n_0^t\le2^{2R^2}
       \le2^{4R^2-4R}\le N^{2R-2}=q/N^2
       \qquad(0\le j<t).
\]
For the derived label, the retained literal constructor bound gives
\[
 J\le N^{3\ell},\quad \ell=2^{22},\qquad
 3\ell(t+1)\le2t!=2F.
\]
The latter inequality holds at \(t\ge16\) in the inherited padded
construction. Consequently, for \(0\le j<t+1\),
\[
 \binom Jj\le\max(1,J^j)\le N^{3\ell t}
       \le N^{2F}\le N^{2R-2}=q/N^2.
\]
No group has more than five members, and \(q/N^2\) is an integer.
Therefore \(5(q/N^2)\le q-1\). Lemma~\ref{lem:digits} constructs all
four new slacks without enlarging the radix or changing the CRT modulus.

\section{One natural row authenticates the Pell indices}
\label{sec:cap}

For \(A\ge2\), define the Pell sequences by
\[
 X_0=1,\ Y_0=0,\qquad
 X_{j+1}=AX_j+(A^2-1)Y_j,\quad Y_{j+1}=X_j+AY_j.
\]
Natural norm-one solutions at discriminant \(A^2-1\) have ordinates
\(Y_j(A)\). The elementary recurrence and multiplication law give
\[
 Y_j\equiv j\pmod{A-1},\quad Y_j\ge j,\quad
 X_j>(A-1)Y_j\ (j\ge1),\quad
 Y_{j+1}>(2A-1)Y_j,
\]
and \(Y_{2j}=2X_jY_j\). Classification can be proved by descent using
the inverse of the smallest positive unit \(A+\sqrt{A^2-1}\).

The inherited frame gives, before interpreting any binomial,
\[
 W\ge3,\quad A-1>W^3,\quad 2\le H<W,\quad
 1\le k_i<W,\quad r=14\le W-2.
\]
A complete retained indexed certificate fixes \(V=Y_H(A)\), where
the semantic anchor is \(H=J\). The other pairs are represented by
consecutive ordinates satisfying
\(Y_i^2+Z_i^2-2AY_iZ_i=1\); retained multiplication gates keep the
literal residuals quadratic. Their output conditions choose the required
orientation, as proved in the inherited source.

\begin{theorem}[Fused natural Pell cap]\label{thm:cap}
Under the frame assumptions above, suppose every \(Y_i\) is a natural
Pell ordinate and \(Q\in\N\) satisfies
\begin{equation}\label{eq:cap}
 \boxed{V=H+\sum_{i=1}^{r}(Y_i-k_i)^2+(A-1)Q.}
\end{equation}
Then \(Y_i=Y_{k_i}(A)\) for all \(i\). Conversely, canonical pairs
extend to a natural \(Q\) whenever \(H\ge2\max_i k_i+1\).
\end{theorem}
\begin{proof}
Write \(Y_i=Y_{j_i}(A)\). Equation~\eqref{eq:cap} gives
\(|Y_i-k_i|\le\sqrt V\). Since \(H\ge2\),
\(V\ge2A>2W\) and \(V>4\). Hence
\[
 Y_i\le k_i+\sqrt V<W+\sqrt V<V.
\]
Strict monotonicity implies \(j_i<H<W\). Reduction modulo \(A-1\)
now shows that \(A-1\) divides
\(E_0=\sum_i(j_i-k_i)^2\). But
\[
 0\le E_0\le r(W-1)^2<W^3<A-1.
\]
Thus \(E_0=0\), proving every requested index.

For the converse put \(m=\max_i k_i\), and take canonical pairs.
Then \(E_1=\sum_i(Y_{k_i}-k_i)^2\le rY_m^2\). The doubled-gap
assumption and Pell growth give
\begin{align*}
 V&\ge Y_{2m+1}>(2A-1)Y_{2m}\\
  &=2(2A-1)X_mY_m
   >2(2A-1)(A-1)Y_m^2.
\end{align*}
Since \(A>W\) and \(r,H<W\), this is greater than
\(rY_m^2+H\). Also \(V-H\) and every \(Y_{k_i}-k_i\) are divisible
by \(A-1\). Therefore \(Q=(V-H-E_1)/(A-1)\) is a nonnegative
integer and satisfies~\eqref{eq:cap}.
\end{proof}

\subsection{The semantic anchor already has the required gap}

The inherited constructive choice gives \(N^2<a\le2N^2\). Hence
\[
 H=J\ge2a^2>2N^4.
\]
Every other CRT seed is at most \(N^2\); the five non-CRT indices
\(t,K,y+1,x+1,F^2+t\) are at most \(N\). Thus
\(m\le2N^2\). For \(N\ge2\), \(2N^4\ge4N^2+1\), so the required
\(H\ge2m+1\) follows with no new coordinate or inequality.

This gap is used only in constructive completeness. Arbitrary zeros need
not satisfy the specific CRT representative choice \(a>N^2\); the
soundness half of Theorem~\ref{thm:cap} does not assume it.

\subsection{The two witness directions in their proper order}

From a new natural zero, first recover the raw frame bounds. The fused
row supplies \(H\le V\), so the retained full anchor criterion applies.
Theorem~\ref{thm:cap} then identifies all fourteen seed ordinates.
Restore the old congruence quotient by
\[
 Q_{\mathrm{old}}=\frac{\sum_i(Y_i-k_i)^2}{A-1}.
\]
It is integral and nonnegative. Restore the old aggregate slack by
\(s=Y_H-H-1-\sum_iY_i\). Since \(k_i<H\) and \(A>W\), growth gives
\(Y_H-rY_{H-1}>H\), so this slack is natural as well. Finally restore
every individual digit slack by Lemma~\ref{lem:digits}. One has now
obtained a zero of the entire inherited graph; its existing soundness
theorem produces the finite refutation.

For completeness, start from the derivation, choose the inherited padded
CRT witnesses and the same lift \(a\mapsto a+N^2\), and compute the
true arithmetic quantities. Section~\ref{sec:direct} supplies the four
digit slacks. Choose the independently bounded Pell parameter in the
retained soluble common-root congruence class. Canonical seed pairs and
the anchor at \(J\), together with Theorem~\ref{thm:cap}, supply \(Q\).
All remaining witnesses are unchanged. This is an equivalence of
\emph{existence}; it is not a claimed bijection on arbitrary old natural
witnesses.

\subsection{Literal contraction to 226 variables}

The actual source transformation starts with 344 natural coordinates and
293 quadratic rows. Removing obsolete cap sum gates, fusing the cap, and
replacing eleven digit bounds by four leaves 332 coordinates and 281 rows.
A checked monic projection removes 106 further coordinates and rows,
giving 226 coordinates and 175 quadratic residuals. The sparse projection
removes 105 and gives 227 coordinates and 176 rows.

Every deleted coordinate has an explicit polynomial restoration certified
nonnegative on the natural domain by nonnegative monomials or rational
sum-of-squares certificates. The final equations are exactly
\[
 P_{226}(\mathbf z)=\sum_{j=0}^{174}Q_j(\mathbf z)^2=0,
 \qquad
 P_{227}(\mathbf z)=\sum_{j=0}^{175}\widetilde Q_j(\mathbf z)^2=0.
\]
All residuals and every collected coefficient are supplied. The principal
residual list is printed in Appendix~\ref{app:residuals}. There are no
variable exponents, factorial symbols, bounded universal quantifiers, or
free coefficient parameters in these literal polynomials.

An optional integer-domain companion substitutes
\[
 z_i=a_i^2+b_i^2+c_i(c_i+1),\qquad 0\le i<226.
\]
This expression is nonnegative at every integer triple and represents
every natural number: apply the three-square theorem to \(4z_i+1\),
which is never of the excluded form \(4^e(8h+7)\), then place the unique
odd square in the third position. The result is a specified degree-eight
polynomial in 678 unrestricted integers. It is a domain conversion, not
an improvement on the specialized signed quartics.

\section{Ten-positive-variable equations of degrees 323 and 318}
\label{sec:compact}

This smaller contribution retains package 133's complete four-coefficient
tuple and native calibration. It is separate from the direct C222 route.
The only variables are
\[
 J,F,\alpha,i,T,w,Y,n,\eta,\zeta\in\Z_{>0}.
\]
The fixed positive integers \(\beta,K_0,M_C,M_F\) are reconstructed by
the finite 943-tile recipe retained in the compact component and printed
in Appendix~\ref{app:constants}. In particular \(\beta\ge31\) is odd
and \(0<M_C<\beta<M_F<2\beta\).

Define the following polynomial abbreviations:
\begin{align*}
 q&=\beta J+1, & h&=q-1, & C&=q-F-\alpha,\\
 X&=wq, & \mu&=q^2+1, & k&=\eta+\zeta,\\
 c&=kY+\eta, & A&=\mu Y(X+1)+1, & \Delta&=A^2-1,\\
 H&=4A-5, & V_0&=X\mu^2Y^2, & E&=1+V_0(V_0+1)k^2,\\
 U&=ic^2, & a&=\Delta U^2E, & F_0&=1+a,\\
 G&=EF_0, & F_s&=1+4aF_0, & Q&=\Delta(F_s-1).
\end{align*}
Also put
\begin{align*}
 R&=[h(q-F)+\alpha](q^2-1)+(M_C+qM_F)J,\\
 e&=k-R-1, & B&=(K_0+w)C-F,\\
 b&=c(T+\alpha)-R, & V&=F_s b-c,\\
 K_+&=X(3c-2X)+2, & P&=q^3h,\qquad M=P\mu^2H.
\end{align*}
The new lower center is
\[
 d=(A-1)UE,
\]
and the carrier is
\begin{align*}
 Z={}&P\mu^2\{H[n+V(d-bB)]+K_+\}\\
     &+H\{\mu^2(hY-q^3B)-Pe\}.
\end{align*}
The single ordinary polynomial equation is
\begin{equation}\label{eq:compact}
 \boxed{Q(Z^2-M^2GV^2)+(M^2G-Z^2)=0.}
\end{equation}
No symbol in these abbreviations adds a quantified unknown.

\subsection{The shear preserves positive solvability}

The 85-operation parent uses the carrier center \(n-bVB\). Let its
polynomial be \(N_{85}\), and let~\eqref{eq:compact} be \(N_{87}\).
The literal ring identity is
\[
 N_{87}(\text{other variables},n)
   =N_{85}(\text{other variables},n+dV).
\]
This alone does not prove positive-domain equivalence: \(V\) must first
be oriented. The parent's pretransport arguments apply to any rational
carrier. From a positive zero they recover positive integers \(g,y\)
with
\[
 G=g^2,\qquad QV^2-(Q-1)y^2=1,
\]
and prove \(b>c\), hence \(V>0\), before using the literal center.
These unchanged arguments use the identities
\[
 F_s=(1+2a)^2,\qquad Q=4\Delta^2U^2G,
\]
which make \(G\) coprime to \((Q-1)(QV^2-1)\). Valuations recover
the square factors and the retained adjacent-Pell descent recovers
integrality. Their full ordinary proof is retained unchanged.
Since \(d>0\), the new zero therefore gives a positive old coordinate
\(n+dV\), so inherited soundness applies.

For completeness use the canonical old witnesses. In the parent's proof
notation let \(s=Y/q^3>0\), \(\kappa=K_+/H>0\), \(\rho=e/\mu^2\ge0\),
\(t_0=B/(q-1)\ge0\), and \(W_0=1+(q-1)bV\). These are integers on
those witnesses, with
\[
 n_{85}=gy+W_0t_0+\rho-\kappa-s.
\]
The exact identity
\begin{equation}\label{eq:centerdefect}
 G-d^2=E+2(A-1)U^2E^2>0
\end{equation}
gives the integer bound \(g-d\ge1\). The canonical inequalities give
\(y>V>2c\) and \(\kappa<c\), \(s<c\). Thus
\begin{align*}
 n_{85}-dV
   &=(g-d)V+g(y-V)+W_0t_0+\rho-\kappa-s\\
   &>V-\kappa-s>0.
\end{align*}
This constructs the new positive coordinate and proves the converse.

\subsection{Exact arithmetic cost and degree}

Both \(A-1\) and \(UE\) are already live in the parent circuit.
Replace its three gates for \(bV,bVB,n-bVB\) by five gates for
\(d,bB,d-bB,V(d-bB),n+V(d-bB)\). The exact count is therefore
87: 52 multiplications, 20 additions, and 15 subtractions. Fixed scalar
products and squares each count as a multiplication.

Let \(Z_0\) be the old uncentered numerator, so \(Z=MdV+Z_0\). Then
the output equals
\begin{align*}
 &-M^2\{Q(G-d^2)+d^2\}V^2
   +2(Q-1)MdVZ_0\\
 &\hspace{30mm}+(Q-1)Z_0^2+M^2G.
\end{align*}
The prefix degrees are
\[
 \deg Q=86,\quad\deg M=13,\quad\deg G=56,\quad
 \deg V=80,\quad\deg d=28,\quad\deg Z_0\le99.
\]
Equation~\eqref{eq:centerdefect} has exact degree 51. The four displayed
terms have degree at most 323, 306, 284, and 82, respectively; the first
has a nonzero highest homogeneous part. With \(q_0=\beta J\), it is
\[
 -2^{11}q_0^{153}w^{33}Y^{77}i^{14}
           (\eta+\zeta)^{44}(q_0-F)^2.
\]
For example the coefficient of
\(J^{153}w^{33}Y^{77}i^{14}\eta^{44}F^2\) is
\(-2^{11}\beta^{153}\ne0\). The literal polynomial therefore has
exact total degree 323, not merely a syntactic upper bound.

This gives a modest new operations/degree tradeoff between the existing
85/328 and 89/288 points. It does not dominate either of them. The
84/436 parent remains the cheaper evaluator among this small set.

\subsection{A sharper center gives 88 operations and degree 318}

There is a second valid choice, \(d_\sharp=AUE\), in exactly the same
carrier and output formula. Compute it by adding the already live
\(UE\) to \((A-1)UE\), so it costs one additional addition. The count
is 88 operations: 52 multiplications, 21 additions, and 15 subtractions.
The shear and new-to-old positivity argument are unchanged.

The lower-center inverse bound \(g>d\) does not apply to this choice.
Instead the inherited canonical source witnesses have
\(X=2^R\), \(R>3q\), \(w=X/q\), \(C\ge1\), and \(F<q\).
For \(q\ge32\), this gives \(w>q>F\), and therefore
\(B=(K_0+w)C-F\ge1\). The retained pretransport bound is \(b>4g\).
Also
\[
 4G-d_\sharp^2=4E+(3A^2-4)U^2E^2>0,
\]
so \(d_\sharp<2g\) and \(bB-d_\sharp>0\). Rewriting the same old
canonical witness gives
\[
 n_{85}-d_\sharp V
   =gy+t_0+\rho-\kappa-s+V(bB-d_\sharp)
   >V-\kappa-s>0.
\]
Thus constructive positive completeness also holds for the sharper
center. This uses actual calibrated witnesses, not an assumption that
the old carrier coordinate can be freely enlarged.

Its exact defect is \(G-d_\sharp^2=E-U^2E^2\), of degree 46.
The four degree bounds in the preceding output identity become
318, 306, 284, and 82. The highest homogeneous part is
\[
 2^{10}(\beta J)^{150}w^{32}Y^{76}i^{14}
             (\eta+\zeta)^{44}(\beta J-F)^2,
\]
which is nonzero. Hence the total degree is exactly 318. The full
component proof further establishes that both shears are bijections
on the respective positive zero sets, including either radical-sign
branch; the existential directions above already suffice for
Theorem~\ref{thm:main}.

\section{Calibration, proof scope, and reproducibility}
\label{sec:calibration}

\subsection{What is proved here and what is imported}

There are three distinct proof layers.
\begin{enumerate}
\item \textbf{New unconditional finite algebra.} The graph theorem proves
an explicit integer zero-set bijection for any checked source system.
The local equality lemma and the exact graph identities do not depend
on a consistency assumption or a proof calculus.
\item \textbf{New unbounded arithmetic.} The shared digit-margin proof,
the natural fused cap, their witness order, and the compact lower-center
argument are ordinary mathematical proofs. Their essential unbounded
claims are not inferred from finite tests.
\item \textbf{Inherited consistency calibration.} The signed324 source
comes from package 131's direct C222 arithmetic theorem; signed442 comes
from package 125. The natural route imports package 129's semantic graph
and its calibration. The named logical interfaces are
\path{KEndpoint.c222_iff_zf_refutation} and
\path{ConstructibleChoice.zf_refutation_iff_zfc_refutation}.
The compact branch retains package 133's separate native/K215 calibration.
\end{enumerate}

The direct source represents a finite syntactic derivation in the
numerical calculus C222. Its retained soundness and completeness use
factorial fixed points, binomial digits, CRT representatives and Horn
derivation decoding. The arithmetic system has a zero exactly when the
calculus has a derivation; the two named logical interfaces then connect
that derivation to a finite ZFC refutation. These are the precise
mathematical dependencies of Theorem~\ref{thm:main}.

No new complete Lean/Rocq development or encoded object-language ZFC proof
term is supplied. The existing formal dependency closures were not all
rebuilt during this work. This is an ordinary proof formalizable in ZFC,
with independently checked finite polynomial artifacts. As in standard
mathematical reuse, importing a named source theorem is distinct from
claiming to have re-proved its entire foundation. The archive retains
source identities, relevant proof texts and explicit boundaries.

\subsection{What the exact checks establish}

The graph verifier is independent of the producer. It checks triangular
normalization, exact restoration of every original residual, complete
alias and unit partitions, all logical block sums, the internal spanning
trees, and the entire collected coefficient list. It also checks the
symbolic canonical lift in~\eqref{eq:lift}. Deliberate mutations of source
rows, incidence signs, alias labels, free flows and the final constant
are rejected. Fresh rebuilds reproduce the mathematical JSON artifacts
byte for byte.

The direct checker independently reconstructs the changed raw rows,
including the four pooled digit bounds and the fused cap. It retains
the unchanged factorial output bound, checks the entire affine export,
replays every monic projection with its nonnegativity certificate, and
compares every expanded quartic coefficient. A separate reviewer checked
the unbounded digit estimates, anchor gap and witness order.

The compact checker reconstructs all four fixed integers from the finite
printed tile recipe, checks their hashes and bit lengths, verifies the
literal circuit modification and cost, and checks the degree certificate.
Its finite-field coefficient controls supplement the displayed symbolic
upper bound; they are not the proof of positive solvability.

Integrity hashes identify delivered bytes; they are not theorem-prover
certificates or claims of an external signature. ``Independent review''
here means a separately tasked AI-assisted review and separately written
checkers, not an external human endorsement.

\subsection{How to use the archive}

The full deliverable is self-contained for the supported finite
reconstruction and coefficient checks. Python 3.12 is the tested runtime;
the implementation uses the standard library. From the archive root:
\begin{lstlisting}[language=bash]
python3 -B verify_all.py --rebuild
python3 -B -O verify_all.py
./build_report.sh
\end{lstlisting}
The PDF rebuild uses a normal pdfLaTeX installation. PDF bytes may differ
across TeX installations; byte equality is required for mathematical data,
not for cross-version PDF metadata. The top-level manifest records the
delivered files. Component READMEs identify smaller commands and exact
artifact formats.

Every graph monomial is an explicit product of indexed coordinates with
coefficient \(1\) or \(-1\). Every direct residual is an explicit integer
quadratic. Every compact register is a polynomial abbreviation in the ten
declared unknowns and four fixed reconstructed integers. No delivered
equation asks the reader to select an unspecified polynomial by DPRM.

\section{Further research questions}

\textbf{Joint graph layout and alias optimization.}
The search first selects a graph and then identifies compatible aliases.
Can a joint optimization of edge incidence, opposite-occurrence matching
and alias classes reduce the 951-variable endpoint substantially? The
complete coefficient certificate permits much more freedom than a global
Sidon or four-cycle restriction. A useful objective should record both
dimension and support, since fewer hubs can increase monomial count.

\textbf{Equality penalties and the actual energy.}
The four-term quartic is locally optimal in the narrow class proved here.
It need not be optimal as part of a larger nonnegative polynomial.
Can cancellations with the graph energy eliminate the remaining dummy
coordinates while preserving height one? Any proposal must certify the
sum's nonnegativity and its zero set, not just its coefficient list.

\textbf{Pool more arithmetic inequalities.}
The direct gain came from proving that several canonical remainders
jointly fit below one denominator. Other source bounds may have comparable
unused margins. Which can be pooled without a circular dependence on the
factorial or binomial interpretation? A sufficient condition must be proved
for canonical witnesses, and a literal restoration map must handle every
new zero.

\textbf{Transfer the pooled digits to the signed source.}
The current 226-variable reduction is over \(\N\). Signed systems recover
remainder signs from product guards. Sharing their bounds may remove more
variables, but requires a fresh sign argument; the natural-domain proof
does not automatically transfer. A smaller signed source would also
improve the height-one compiler's output.

\textbf{Improve the compact center.}
The center \((A-1)UE\) cancels five leading degrees at a cost of two gates.
Can another polynomial lower center reduce the degree more while preserving
the explicit positive reconstruction? A formal coordinate shear is not
enough: positivity of the transformed witness must be proved in both
directions and in the correct order.

\textbf{Close the formal calibration chain.}
The graph transformation is a small, general target for proof-assistant
formalization: finite maps, integer polynomial identities, positive
semidefinite energies, and unimodular tree elimination. The direct
arithmetic bridge is more substantial. A future formal release should
make the exact internal-ZFC endpoint explicit rather than treating
finite script checks as a substitute.

\textbf{Meaningful lower bounds.}
The only optimality statement here is the local two-variable equality
gadget theorem. The global minimum number of variables, degree, support,
coefficient height or operations for a Con(ZFC) equation remains unknown.
A valuable next step would be a lower bound for a precisely specified
compiler class, not an unsupported claim that one displayed equation is
the simplest possible.

\appendix
\section{Complete principal natural residuals}\label{app:residuals}

The following is the complete 175-row quadratic specification of
\(P_{226}=\sum_jQ_j^2\). The variables are \(z_0,\ldots,z_{225}\in\N\).
Rows are generated directly from the authoritative equation JSON and
checked against it. The collected expansion is also included in the
archive; the row format is substantially shorter to read.
% Generated from direct/equations/natural226/equation.json.
\begingroup\small
\allowdisplaybreaks
\[\begin{aligned}Q_{0}={}&-z_{2}^{2}+z_{3}
\end{aligned}\]
\[\begin{aligned}Q_{1}={}&1-2z_{2}-z_{2}z_{4}+z_{4}+z_{5}-z_{6}
\end{aligned}\]
\[\begin{aligned}Q_{2}={}&1+z_{4}-z_{6}-z_{7}
\end{aligned}\]
\[\begin{aligned}Q_{3}={}&-17-z_{0}-z_{0}z_{2}-z_{0}z_{9}-17z_{2}-17z_{9}
\\&+z_{11}
\end{aligned}\]
\[\begin{aligned}Q_{4}={}&-1+z_{1}-z_{2}-z_{2}z_{8}-2z_{8}-z_{8}z_{9}
\end{aligned}\]
\[\begin{aligned}Q_{5}={}&-1+z_{1}-z_{10}-z_{10}z_{11}
\end{aligned}\]
\[\begin{aligned}Q_{6}={}&2+z_{3}-z_{4}z_{22}+z_{20}-2z_{22}
\end{aligned}\]
\[\begin{aligned}Q_{7}={}&-6-5z_{12}-z_{12}^{2}-2z_{12}z_{13}-5z_{13}-z_{13}^{2}
\\&+z_{23}
\end{aligned}\]
\[\begin{aligned}Q_{8}={}&-72-68z_{13}-16z_{13}^{2}-16z_{13}z_{23}-34z_{23}-4z_{23}^{2}
\\&+z_{24}
\end{aligned}\]
\[\begin{aligned}Q_{9}={}&-30-11z_{12}-z_{12}^{2}-4z_{12}z_{24}-22z_{24}-4z_{24}^{2}
\\&+z_{25}
\end{aligned}\]
\[\begin{aligned}Q_{10}={}&-380-312z_{24}-64z_{24}^{2}-32z_{24}z_{25}-78z_{25}-4z_{25}^{2}
\\&+z_{26}
\end{aligned}\]
\[\begin{aligned}Q_{11}={}&-6-5z_{13}-z_{13}^{2}-2z_{13}z_{14}-5z_{14}-z_{14}^{2}
\\&+z_{27}
\end{aligned}\]
\[\begin{aligned}Q_{12}={}&-6-5z_{12}-z_{12}^{2}-2z_{12}z_{14}-5z_{14}-z_{14}^{2}
\\&+z_{28}
\end{aligned}\]
\[\begin{aligned}Q_{13}={}&-414-22z_{12}-2z_{12}^{2}-16z_{12}z_{14}-8z_{12}z_{27}-200z_{13}
\\&-32z_{13}^{2}-64z_{13}z_{14}-32z_{13}z_{23}-32z_{13}z_{28}-320z_{14}-64z_{14}^{2}
\\&-32z_{14}z_{23}-32z_{14}z_{27}-32z_{14}z_{28}-100z_{23}-8z_{23}^{2}-16z_{23}z_{28}
\\&-52z_{27}-8z_{27}^{2}-108z_{28}-8z_{28}^{2}+z_{29}
\end{aligned}\]
\[\begin{aligned}Q_{14}={}&-1344-800z_{13}-128z_{13}^{2}-256z_{13}z_{14}-128z_{13}z_{23}-128z_{13}z_{28}
\\&-864z_{14}-128z_{14}^{2}-128z_{14}z_{23}-128z_{14}z_{28}-400z_{23}-32z_{23}^{2}
\\&-64z_{23}z_{28}-432z_{28}-32z_{28}^{2}-2z_{29}-2z_{29}^{2}+z_{30}
\end{aligned}\]
\[\begin{aligned}Q_{15}={}&-247-168z_{12}-32z_{12}^{2}-64z_{12}z_{13}-32z_{12}z_{27}-32z_{12}z_{28}
\\&-188z_{13}-32z_{13}^{2}-32z_{13}z_{27}-32z_{13}z_{28}-2z_{23}-92z_{27}
\\&-8z_{27}^{2}-16z_{27}z_{28}-84z_{28}-8z_{28}^{2}+z_{31}
\end{aligned}\]
\[\begin{aligned}Q_{16}={}&-960-672z_{12}-128z_{12}^{2}-256z_{12}z_{13}-128z_{12}z_{27}-128z_{12}z_{28}
\\&-736z_{13}-128z_{13}^{2}-128z_{13}z_{27}-128z_{13}z_{28}-368z_{27}-32z_{27}^{2}
\\&-64z_{27}z_{28}-336z_{28}-32z_{28}^{2}-2z_{31}-2z_{31}^{2}+z_{32}
\end{aligned}\]
\[\begin{aligned}Q_{17}={}&-276-30z_{12}-2z_{12}^{2}-16z_{12}z_{14}-8z_{12}z_{27}-30z_{13}
\\&-2z_{13}^{2}-16z_{13}z_{14}-8z_{13}z_{28}-272z_{14}-64z_{14}^{2}-32z_{14}z_{27}
\\&-32z_{14}z_{28}-68z_{27}-8z_{27}^{2}-68z_{28}-8z_{28}^{2}+z_{33}
\end{aligned}\]
\[\begin{aligned}Q_{18}={}&-552-120z_{13}-8z_{13}^{2}-64z_{13}z_{14}-32z_{13}z_{28}-544z_{14}
\\&-128z_{14}^{2}-128z_{14}z_{28}-272z_{28}-32z_{28}^{2}-2z_{33}-2z_{33}^{2}
\\&+z_{34}
\end{aligned}\]
\[\begin{aligned}Q_{19}={}&-247-4z_{12}-352z_{13}-128z_{13}^{2}-64z_{13}z_{23}-64z_{13}z_{27}
\\&-92z_{23}-8z_{23}^{2}-16z_{23}z_{27}-84z_{27}-8z_{27}^{2}-2z_{28}
\\&+z_{35}
\end{aligned}\]
\[\begin{aligned}Q_{20}={}&-960-1408z_{13}-512z_{13}^{2}-256z_{13}z_{23}-256z_{13}z_{27}-368z_{23}
\\&-32z_{23}^{2}-64z_{23}z_{27}-336z_{27}-32z_{27}^{2}-2z_{35}-2z_{35}^{2}
\\&+z_{36}
\end{aligned}\]
\[\begin{aligned}Q_{21}={}&-56-60z_{13}-16z_{13}^{2}-16z_{13}z_{23}-30z_{23}-4z_{23}^{2}
\\&+z_{37}
\end{aligned}\]
\[\begin{aligned}Q_{22}={}&-30-11z_{12}-z_{12}^{2}-4z_{12}z_{37}-22z_{37}-4z_{37}^{2}
\\&+z_{38}
\end{aligned}\]
\[\begin{aligned}Q_{23}={}&-380-312z_{37}-64z_{37}^{2}-32z_{37}z_{38}-78z_{38}-4z_{38}^{2}
\\&+z_{39}
\end{aligned}\]
\[\begin{aligned}Q_{24}={}&-30-11z_{13}-z_{13}^{2}-4z_{13}z_{39}-22z_{39}-4z_{39}^{2}
\\&+z_{40}
\end{aligned}\]
\[\begin{aligned}Q_{25}={}&-349-420z_{13}-128z_{13}^{2}-128z_{13}z_{23}-208z_{23}-32z_{23}^{2}
\\&-2z_{27}+z_{41}
\end{aligned}\]
\[\begin{aligned}Q_{26}={}&-1368-1664z_{13}-512z_{13}^{2}-512z_{13}z_{23}-832z_{23}-128z_{23}^{2}
\\&-2z_{41}-2z_{41}^{2}+z_{42}
\end{aligned}\]
\[\begin{aligned}Q_{27}={}&-56-60z_{12}-16z_{12}^{2}-16z_{12}z_{28}-8z_{12}z_{42}-30z_{28}
\\&-4z_{28}^{2}-4z_{28}z_{42}-15z_{42}-z_{42}^{2}+z_{43}
\end{aligned}\]
\[\begin{aligned}Q_{28}={}&-6-5z_{14}-z_{14}^{2}-2z_{14}z_{15}-5z_{15}-z_{15}^{2}
\\&+z_{44}
\end{aligned}\]
\[\begin{aligned}Q_{29}={}&-224-26z_{12}-2z_{12}^{2}-16z_{12}z_{15}-8z_{12}z_{44}-288z_{14}
\\&-128z_{14}^{2}-64z_{14}z_{27}-64z_{14}z_{28}-88z_{15}-32z_{15}^{2}-32z_{15}z_{44}
\\&-76z_{27}-8z_{27}^{2}-16z_{27}z_{28}-68z_{28}-8z_{28}^{2}-44z_{44}
\\&-8z_{44}^{2}+z_{45}
\end{aligned}\]
\[\begin{aligned}Q_{30}={}&-640-1152z_{14}-512z_{14}^{2}-256z_{14}z_{27}-256z_{14}z_{28}-304z_{27}
\\&-32z_{27}^{2}-64z_{27}z_{28}-272z_{28}-32z_{28}^{2}-2z_{45}-2z_{45}^{2}
\\&+z_{46}
\end{aligned}\]
\[\begin{aligned}Q_{31}={}&-241-168z_{12}-32z_{12}^{2}-64z_{12}z_{13}-32z_{12}z_{27}-32z_{12}z_{28}
\\&-184z_{13}-32z_{13}^{2}-32z_{13}z_{27}-32z_{13}z_{28}-z_{14}-92z_{27}
\\&-8z_{27}^{2}-16z_{27}z_{28}-84z_{28}-8z_{28}^{2}+z_{47}
\end{aligned}\]
\[\begin{aligned}Q_{32}={}&-42-52z_{14}-16z_{14}^{2}-16z_{14}z_{27}-26z_{27}-4z_{27}^{2}
\\&+z_{48}
\end{aligned}\]
\[\begin{aligned}Q_{33}={}&-30-11z_{14}-z_{14}^{2}-4z_{14}z_{48}-22z_{48}-4z_{48}^{2}
\\&+z_{49}
\end{aligned}\]
\[\begin{aligned}Q_{34}={}&-306-35z_{13}-z_{13}^{2}-16z_{13}z_{48}-4z_{13}z_{49}-280z_{48}
\\&-64z_{48}^{2}-32z_{48}z_{49}-70z_{49}-4z_{49}^{2}+z_{50}
\end{aligned}\]
\[\begin{aligned}Q_{35}={}&-4556-4320z_{48}-1024z_{48}^{2}-512z_{48}z_{49}-128z_{48}z_{50}-1080z_{49}
\\&-64z_{49}^{2}-32z_{49}z_{50}-270z_{50}-4z_{50}^{2}+z_{51}
\end{aligned}\]
\[\begin{aligned}Q_{36}={}&-30-11z_{14}-z_{14}^{2}-4z_{14}z_{51}-22z_{51}-4z_{51}^{2}
\\&+z_{52}
\end{aligned}\]
\[\begin{aligned}Q_{37}={}&-306-35z_{13}-z_{13}^{2}-16z_{13}z_{51}-4z_{13}z_{52}-280z_{51}
\\&-64z_{51}^{2}-32z_{51}z_{52}-70z_{52}-4z_{52}^{2}+z_{53}
\end{aligned}\]
\[\begin{aligned}Q_{38}={}&-1028-672z_{12}-128z_{12}^{2}-256z_{12}z_{13}-128z_{12}z_{27}-128z_{12}z_{28}
\\&-736z_{13}-128z_{13}^{2}-128z_{13}z_{27}-128z_{13}z_{28}-368z_{27}-32z_{27}^{2}
\\&-64z_{27}z_{28}-336z_{28}-32z_{28}^{2}-2z_{47}-2z_{47}^{2}-32z_{51}
\\&-8z_{52}-2z_{53}+z_{54}
\end{aligned}\]
\[\begin{aligned}Q_{39}={}&-264-128z_{51}-32z_{52}-8z_{53}-2z_{54}-2z_{54}^{2}
\\&+z_{55}
\end{aligned}\]
\[\begin{aligned}Q_{40}={}&-2-3z_{12}-z_{12}^{2}-2z_{12}z_{55}-3z_{55}-z_{55}^{2}
\\&+z_{56}
\end{aligned}\]
\[\begin{aligned}Q_{41}={}&-12-28z_{55}-16z_{55}^{2}-16z_{55}z_{56}-14z_{56}-4z_{56}^{2}
\\&+z_{57}
\end{aligned}\]
\[\begin{aligned}Q_{42}={}&-30-11z_{14}-z_{14}^{2}-4z_{14}z_{57}-22z_{57}-4z_{57}^{2}
\\&+z_{58}
\end{aligned}\]
\[\begin{aligned}Q_{43}={}&-306-35z_{12}-z_{12}^{2}-16z_{12}z_{57}-4z_{12}z_{58}-280z_{57}
\\&-64z_{57}^{2}-32z_{57}z_{58}-70z_{58}-4z_{58}^{2}+z_{59}
\end{aligned}\]
\[\begin{aligned}Q_{44}={}&-4556-4320z_{57}-1024z_{57}^{2}-512z_{57}z_{58}-128z_{57}z_{59}-1080z_{58}
\\&-64z_{58}^{2}-32z_{58}z_{59}-270z_{59}-4z_{59}^{2}+z_{60}
\end{aligned}\]
\[\begin{aligned}Q_{45}={}&-90-76z_{14}-16z_{14}^{2}-16z_{14}z_{27}-16z_{14}z_{60}-38z_{27}
\\&-4z_{27}^{2}-8z_{27}z_{60}-38z_{60}-4z_{60}^{2}+z_{61}
\end{aligned}\]
\[\begin{aligned}Q_{46}={}&-4556-4320z_{51}-1024z_{51}^{2}-512z_{51}z_{52}-128z_{51}z_{53}-1080z_{52}
\\&-64z_{52}^{2}-32z_{52}z_{53}-270z_{53}-4z_{53}^{2}+z_{62}
\end{aligned}\]
\[\begin{aligned}Q_{47}={}&-682-70z_{13}-2z_{13}^{2}-32z_{13}z_{60}-8z_{13}z_{61}-592z_{60}
\\&-128z_{60}^{2}-64z_{60}z_{61}-148z_{61}-8z_{61}^{2}-2z_{62}+z_{63}
\end{aligned}\]
\[\begin{aligned}Q_{48}={}&-z_{63}-z_{63}^{2}+z_{64}
\end{aligned}\]
\[\begin{aligned}Q_{49}={}&-42-52z_{13}-16z_{13}^{2}-8z_{13}z_{14}-16z_{13}z_{23}-13z_{14}
\\&-z_{14}^{2}-4z_{14}z_{23}-26z_{23}-4z_{23}^{2}+z_{65}
\end{aligned}\]
\[\begin{aligned}Q_{50}={}&-2278-2264z_{13}-800z_{13}^{2}-800z_{13}z_{23}-160z_{13}z_{65}-70z_{14}
\\&-2z_{14}^{2}-32z_{14}z_{62}-8z_{14}z_{64}-1132z_{23}-200z_{23}^{2}-80z_{23}z_{65}
\\&-592z_{62}-128z_{62}^{2}-64z_{62}z_{64}-148z_{64}-8z_{64}^{2}-228z_{65}
\\&-8z_{65}^{2}+z_{66}
\end{aligned}\]
\[\begin{aligned}Q_{51}={}&-2712-280z_{14}-8z_{14}^{2}-128z_{14}z_{62}-32z_{14}z_{64}-2368z_{62}
\\&-512z_{62}^{2}-256z_{62}z_{64}-592z_{64}-32z_{64}^{2}-2z_{66}-2z_{66}^{2}
\\&+z_{67}
\end{aligned}\]
\[\begin{aligned}Q_{52}={}&-2-3z_{67}-z_{67}^{2}+z_{68}
\end{aligned}\]
\[\begin{aligned}Q_{53}={}&-z_{2}-z_{2}z_{17}+z_{69}
\end{aligned}\]
\[\begin{aligned}Q_{54}={}&-z_{16}^{2}+2z_{16}z_{17}-z_{17}^{2}+z_{70}
\end{aligned}\]
\[\begin{aligned}Q_{55}={}&-2z_{26}z_{69}-4z_{69}+z_{71}
\end{aligned}\]
\[\begin{aligned}Q_{56}={}&-z_{16}^{2}+2z_{16}z_{17}+2z_{16}z_{71}-z_{17}^{2}-2z_{17}z_{71}-z_{71}^{2}
\\&+z_{72}
\end{aligned}\]
\[\begin{aligned}Q_{57}={}&-z_{30}z_{69}+z_{73}
\end{aligned}\]
\[\begin{aligned}Q_{58}={}&-z_{16}^{2}+2z_{16}z_{17}+2z_{16}z_{73}-z_{17}^{2}-2z_{17}z_{73}-z_{73}^{2}
\\&+z_{74}
\end{aligned}\]
\[\begin{aligned}Q_{59}={}&-z_{32}z_{69}+z_{75}
\end{aligned}\]
\[\begin{aligned}Q_{60}={}&-z_{16}^{2}+2z_{16}z_{17}+2z_{16}z_{75}-z_{17}^{2}-2z_{17}z_{75}-z_{75}^{2}
\\&+z_{76}
\end{aligned}\]
\[\begin{aligned}Q_{61}={}&-z_{34}z_{69}+z_{77}
\end{aligned}\]
\[\begin{aligned}Q_{62}={}&-z_{16}^{2}+2z_{16}z_{17}+2z_{16}z_{77}-z_{17}^{2}-2z_{17}z_{77}-z_{77}^{2}
\\&+z_{78}
\end{aligned}\]
\[\begin{aligned}Q_{63}={}&-z_{36}z_{69}+z_{79}
\end{aligned}\]
\[\begin{aligned}Q_{64}={}&-z_{16}^{2}+2z_{16}z_{17}+2z_{16}z_{79}-z_{17}^{2}-2z_{17}z_{79}-z_{79}^{2}
\\&+z_{80}
\end{aligned}\]
\[\begin{aligned}Q_{65}={}&-4z_{13}z_{69}-2z_{40}z_{69}-4z_{69}+z_{81}
\end{aligned}\]
\[\begin{aligned}Q_{66}={}&-z_{16}^{2}+2z_{16}z_{17}+2z_{16}z_{81}-z_{17}^{2}-2z_{17}z_{81}-z_{81}^{2}
\\&+z_{82}
\end{aligned}\]
\[\begin{aligned}Q_{67}={}&-4z_{42}z_{69}-2z_{43}z_{69}+z_{83}
\end{aligned}\]
\[\begin{aligned}Q_{68}={}&-z_{16}^{2}+2z_{16}z_{17}+2z_{16}z_{83}-z_{17}^{2}-2z_{17}z_{83}-z_{83}^{2}
\\&+z_{84}
\end{aligned}\]
\[\begin{aligned}Q_{69}={}&-z_{46}z_{69}+z_{85}
\end{aligned}\]
\[\begin{aligned}Q_{70}={}&-z_{16}^{2}+2z_{16}z_{17}+2z_{16}z_{85}-z_{17}^{2}-2z_{17}z_{85}-z_{85}^{2}
\\&+z_{86}
\end{aligned}\]
\[\begin{aligned}Q_{71}={}&-z_{12}^{2}+2z_{12}z_{14}-z_{14}^{2}+z_{87}
\end{aligned}\]
\[\begin{aligned}Q_{72}={}&-z_{13}^{2}+2z_{13}z_{14}-z_{14}^{2}+z_{88}
\end{aligned}\]
\[\begin{aligned}Q_{73}={}&-z_{87}z_{88}+z_{89}
\end{aligned}\]
\[\begin{aligned}Q_{74}={}&-4-4z_{15}-z_{15}^{2}+z_{90}
\end{aligned}\]
\[\begin{aligned}Q_{75}={}&-z_{89}^{2}+2z_{89}z_{90}-z_{90}^{2}+z_{91}
\end{aligned}\]
\[\begin{aligned}Q_{76}={}&-2z_{68}z_{69}-4z_{69}+z_{92}
\end{aligned}\]
\[\begin{aligned}Q_{77}={}&-z_{16}^{2}+2z_{16}z_{17}+2z_{16}z_{92}-z_{17}^{2}-2z_{17}z_{92}-z_{92}^{2}
\\&+z_{93}
\end{aligned}\]
\[\begin{aligned}Q_{78}={}&-z_{2}-z_{2}z_{13}+z_{94}
\end{aligned}\]
\[\begin{aligned}Q_{79}={}&-1-z_{17}-z_{17}z_{94}-z_{94}+z_{95}
\end{aligned}\]
\[\begin{aligned}Q_{80}={}&-z_{2}-z_{2}z_{12}+z_{96}
\end{aligned}\]
\[\begin{aligned}Q_{81}={}&-4z_{2}-4z_{2}z_{13}-2z_{2}z_{23}+z_{97}
\end{aligned}\]
\[\begin{aligned}Q_{82}={}&-1-z_{16}-z_{16}z_{96}-z_{96}+z_{98}
\end{aligned}\]
\[\begin{aligned}Q_{83}={}&-z_{97}z_{98}-z_{98}+z_{99}
\end{aligned}\]
\[\begin{aligned}Q_{84}={}&-z_{95}^{2}+2z_{95}z_{99}-z_{99}^{2}+z_{100}
\end{aligned}\]
\[\begin{aligned}Q_{85}={}&-6z_{2}-4z_{2}z_{13}-2z_{2}z_{23}+z_{101}
\end{aligned}\]
\[\begin{aligned}Q_{86}={}&-1-z_{17}-z_{17}z_{101}-z_{101}+z_{102}
\end{aligned}\]
\[\begin{aligned}Q_{87}={}&-1-z_{16}-z_{16}z_{94}-z_{94}+z_{103}
\end{aligned}\]
\[\begin{aligned}Q_{88}={}&-z_{102}^{2}+2z_{102}z_{103}-z_{103}^{2}+z_{104}
\end{aligned}\]
\[\begin{aligned}Q_{89}={}&-z_{70}z_{72}+z_{105}
\end{aligned}\]
\[\begin{aligned}Q_{90}={}&-z_{74}z_{105}+z_{106}
\end{aligned}\]
\[\begin{aligned}Q_{91}={}&-z_{76}z_{106}+z_{107}
\end{aligned}\]
\[\begin{aligned}Q_{92}={}&-z_{78}z_{107}+z_{108}
\end{aligned}\]
\[\begin{aligned}Q_{93}={}&-z_{80}z_{108}+z_{109}
\end{aligned}\]
\[\begin{aligned}Q_{94}={}&-z_{82}z_{109}+z_{110}
\end{aligned}\]
\[\begin{aligned}Q_{95}={}&-z_{84}z_{110}+z_{111}
\end{aligned}\]
\[\begin{aligned}Q_{96}={}&-z_{86}z_{111}+z_{112}
\end{aligned}\]
\[\begin{aligned}Q_{97}={}&-z_{91}z_{112}-z_{93}z_{112}+z_{113}
\end{aligned}\]
\[\begin{aligned}Q_{98}={}&-z_{100}z_{113}+z_{114}
\end{aligned}\]
\[\begin{aligned}Q_{99}={}&-2-2z_{2}-z_{2}z_{20}-2z_{9}-z_{9}z_{20}-z_{20}
\\&+z_{115}
\end{aligned}\]
\[\begin{aligned}Q_{100}={}&-1-z_{18}-z_{18}z_{115}-z_{115}+z_{116}
\end{aligned}\]
\[\begin{aligned}Q_{101}={}&-2-z_{2}-z_{2}z_{19}-z_{9}-z_{9}z_{19}-2z_{19}
\\&-z_{19}z_{115}-z_{115}+z_{117}
\end{aligned}\]
\[\begin{aligned}Q_{102}={}&2-4z_{1}+2z_{1}^{2}-2z_{1}z_{16}-2z_{1}z_{17}-2z_{1}z_{116}
\\&-2z_{1}z_{117}+2z_{16}+z_{16}^{2}+2z_{16}z_{116}+2z_{17}+z_{17}^{2}
\\&+2z_{17}z_{117}-z_{21}z_{118}+z_{104}z_{114}+2z_{116}+z_{116}^{2}+2z_{117}
\\&+z_{117}^{2}
\end{aligned}\]
\[\begin{aligned}Q_{103}={}&-9-z_{12}-z_{13}-z_{14}-z_{15}-z_{16}
\\&-z_{17}-z_{18}-z_{19}-z_{20}+10z_{21}^{2}-z_{119}
\end{aligned}\]
\[\begin{aligned}Q_{104}={}&-2-z_{4}-z_{4}z_{169}-z_{4}z_{170}-z_{4}z_{221}-3z_{169}
\\&-z_{169}z_{171}-3z_{170}-z_{170}z_{171}-z_{171}-z_{171}z_{221}+z_{176}
\\&+z_{177}+z_{183}+z_{188}+z_{193}-3z_{221}+z_{224}
\end{aligned}\]
\[\begin{aligned}Q_{105}={}&-132-4z_{0}-4z_{0}z_{122}-3z_{1}-3z_{1}z_{122}-5z_{2}
\\&-5z_{2}z_{122}-2z_{3}-2z_{3}z_{122}-2z_{4}-2z_{4}z_{122}-z_{5}
\\&-z_{5}z_{122}-3z_{9}-3z_{9}z_{122}-2z_{12}-2z_{12}z_{122}-14z_{13}
\\&-14z_{13}z_{122}-z_{14}-z_{14}z_{122}-z_{15}-z_{15}z_{122}-z_{16}
\\&-z_{16}z_{122}-z_{17}-z_{17}z_{122}-z_{18}-z_{18}z_{122}-z_{19}
\\&-z_{19}z_{122}-z_{20}-z_{20}z_{122}-z_{21}-z_{21}z_{122}-4z_{23}
\\&-4z_{23}z_{122}-2z_{26}-2z_{26}z_{122}-z_{30}-z_{30}z_{122}-z_{32}
\\&-z_{32}z_{122}-z_{34}-z_{34}z_{122}-z_{36}-z_{36}z_{122}-2z_{40}
\\&-2z_{40}z_{122}-4z_{42}-4z_{42}z_{122}-2z_{43}-2z_{43}z_{122}-z_{46}
\\&-z_{46}z_{122}-2z_{68}-2z_{68}z_{122}-132z_{122}-z_{122}z_{169}-z_{122}z_{170}
\\&-z_{122}z_{171}-z_{122}z_{198}-z_{122}z_{199}-z_{122}z_{205}-z_{122}z_{210}-z_{122}z_{211}
\\&-z_{122}z_{216}-z_{122}z_{221}+z_{123}-z_{169}-z_{170}-z_{171}
\\&-z_{198}-z_{199}-z_{205}-z_{210}-z_{211}-z_{216}
\\&-z_{221}
\end{aligned}\]
\[\begin{aligned}Q_{106}={}&-4z_{0}z_{123}-3z_{1}z_{123}-5z_{2}z_{123}-2z_{3}z_{123}-2z_{4}z_{123}-z_{5}z_{123}
\\&-3z_{9}z_{123}-2z_{12}z_{123}-14z_{13}z_{123}-z_{14}z_{123}-z_{15}z_{123}-z_{16}z_{123}
\\&-z_{17}z_{123}-z_{18}z_{123}-z_{19}z_{123}-z_{20}z_{123}-z_{21}z_{123}-4z_{23}z_{123}
\\&-2z_{26}z_{123}-z_{30}z_{123}-z_{32}z_{123}-z_{34}z_{123}-z_{36}z_{123}-2z_{40}z_{123}
\\&-4z_{42}z_{123}-2z_{43}z_{123}-z_{46}z_{123}-2z_{68}z_{123}-133z_{123}-z_{123}z_{169}
\\&-z_{123}z_{170}-z_{123}z_{171}-z_{123}z_{198}-z_{123}z_{199}-z_{123}z_{205}-z_{123}z_{210}
\\&-z_{123}z_{211}-z_{123}z_{216}-z_{123}z_{221}+z_{124}
\end{aligned}\]
\[\begin{aligned}Q_{107}={}&-1+z_{120}^{2}+z_{123}^{2}-z_{124}^{2}
\end{aligned}\]
\[\begin{aligned}Q_{108}={}&-17-8z_{4}-z_{4}^{2}+2z_{4}z_{120}+2z_{4}z_{121}-2z_{4}z_{171}
\\&+8z_{120}+2z_{120}z_{171}+8z_{121}+2z_{121}z_{171}-z_{125}z_{126}-8z_{171}
\\&-z_{171}^{2}
\end{aligned}\]
\[\begin{aligned}Q_{109}={}&-z_{21}-z_{21}z_{210}-z_{21}z_{211}-z_{21}z_{216}+z_{126}
\end{aligned}\]
\[\begin{aligned}Q_{110}={}&-4z_{128}-4z_{128}z_{133}+z_{134}
\end{aligned}\]
\[\begin{aligned}Q_{111}={}&-z_{120}z_{128}-z_{121}z_{128}+z_{135}
\end{aligned}\]
\[\begin{aligned}Q_{112}={}&-1+z_{127}^{2}+z_{128}^{2}-z_{135}^{2}
\end{aligned}\]
\[\begin{aligned}Q_{113}={}&-z_{120}z_{130}-z_{121}z_{130}+z_{136}
\end{aligned}\]
\[\begin{aligned}Q_{114}={}&-1+z_{129}^{2}+z_{130}^{2}-z_{136}^{2}
\end{aligned}\]
\[\begin{aligned}Q_{115}={}&-z_{132}-z_{132}z_{134}+z_{137}
\end{aligned}\]
\[\begin{aligned}Q_{116}={}&-1+z_{131}^{2}+z_{132}^{2}-z_{137}^{2}
\end{aligned}\]
\[\begin{aligned}Q_{117}={}&-z_{128}^{2}+z_{139}
\end{aligned}\]
\[\begin{aligned}Q_{118}={}&z_{130}-z_{138}z_{139}-z_{139}
\end{aligned}\]
\[\begin{aligned}Q_{119}={}&1-z_{120}-z_{121}-z_{129}z_{140}+z_{134}
\end{aligned}\]
\[\begin{aligned}Q_{120}={}&-z_{127}-z_{129}z_{141}+z_{131}
\end{aligned}\]
\[\begin{aligned}Q_{121}={}&-25-z_{12}-13z_{13}-4z_{23}-2z_{26}-z_{30}
\\&-z_{32}-z_{34}-z_{36}-2z_{40}-4z_{42}-2z_{43}
\\&-z_{46}-2z_{68}-4z_{128}z_{142}+z_{132}
\end{aligned}\]
\[\begin{aligned}Q_{122}={}&-z_{143}z_{144}+z_{145}
\end{aligned}\]
\[\begin{aligned}Q_{123}={}&-1-2z_{120}z_{145}-2z_{121}z_{145}+z_{143}^{2}+z_{144}^{2}
\end{aligned}\]
\[\begin{aligned}Q_{124}={}&-2-2z_{3}-z_{3}^{2}+2z_{3}z_{120}+2z_{3}z_{121}+2z_{120}
\\&+2z_{121}-z_{146}
\end{aligned}\]
\[\begin{aligned}Q_{125}={}&z_{3}z_{143}-z_{5}+z_{143}-z_{144}-z_{146}z_{147}
\end{aligned}\]
\[\begin{aligned}Q_{126}={}&-z_{148}z_{149}+z_{150}
\end{aligned}\]
\[\begin{aligned}Q_{127}={}&-1-2z_{120}z_{150}-2z_{121}z_{150}+z_{148}^{2}+z_{149}^{2}
\end{aligned}\]
\[\begin{aligned}Q_{128}={}&-5-4z_{4}-z_{4}^{2}+2z_{4}z_{120}+2z_{4}z_{121}+4z_{120}
\\&+4z_{121}-z_{151}
\end{aligned}\]
\[\begin{aligned}Q_{129}={}&z_{4}z_{148}-z_{21}+2z_{148}-z_{149}-z_{151}z_{152}
\end{aligned}\]
\[\begin{aligned}Q_{130}={}&-1-z_{21}^{2}+2z_{21}z_{120}+2z_{21}z_{121}-z_{153}
\end{aligned}\]
\[\begin{aligned}Q_{131}={}&-3-z_{4}+z_{21}z_{148}-z_{149}-z_{153}z_{154}-z_{171}
\end{aligned}\]
\[\begin{aligned}Q_{132}={}&-10-6z_{4}-z_{4}^{2}+2z_{4}z_{120}+2z_{4}z_{121}-2z_{4}z_{171}
\\&+6z_{120}+2z_{120}z_{171}+6z_{121}+2z_{121}z_{171}-z_{155}-6z_{171}
\\&-z_{171}^{2}
\end{aligned}\]
\[\begin{aligned}Q_{133}={}&-1+z_{4}z_{143}+3z_{143}+z_{143}z_{171}-z_{144}-z_{155}z_{156}
\\&-z_{169}-z_{170}-z_{221}
\end{aligned}\]
\[\begin{aligned}Q_{134}={}&-z_{157}z_{158}+z_{159}
\end{aligned}\]
\[\begin{aligned}Q_{135}={}&-1-2z_{120}z_{159}-2z_{121}z_{159}+z_{157}^{2}+z_{158}^{2}
\end{aligned}\]
\[\begin{aligned}Q_{136}={}&-1+z_{4}z_{157}-z_{155}z_{160}+3z_{157}+z_{157}z_{171}-z_{158}
\\&-z_{198}-z_{199}-z_{205}
\end{aligned}\]
\[\begin{aligned}Q_{137}={}&-z_{161}z_{162}+z_{163}
\end{aligned}\]
\[\begin{aligned}Q_{138}={}&-1-2z_{120}z_{163}-2z_{121}z_{163}+z_{161}^{2}+z_{162}^{2}
\end{aligned}\]
\[\begin{aligned}Q_{139}={}&-1+z_{4}z_{161}-z_{155}z_{164}+3z_{161}+z_{161}z_{171}-z_{162}
\\&-z_{210}-z_{211}-z_{216}
\end{aligned}\]
\[\begin{aligned}Q_{140}={}&-z_{165}z_{166}+z_{167}
\end{aligned}\]
\[\begin{aligned}Q_{141}={}&-1-2z_{120}z_{167}-2z_{121}z_{167}+z_{165}^{2}+z_{166}^{2}
\end{aligned}\]
\[\begin{aligned}Q_{142}={}&-2-z_{4}+z_{4}z_{165}-z_{4}z_{169}-z_{4}z_{170}-z_{4}z_{221}
\\&+4z_{165}+z_{165}z_{171}-z_{166}-z_{168}-z_{168}z_{176}-z_{168}z_{177}
\\&-z_{168}z_{183}-z_{168}z_{188}-z_{168}z_{193}-z_{168}z_{224}-3z_{169}-2z_{170}
\\&-2z_{221}
\end{aligned}\]
\[\begin{aligned}Q_{143}={}&-z_{172}z_{173}+z_{174}
\end{aligned}\]
\[\begin{aligned}Q_{144}={}&-1-2z_{120}z_{174}-2z_{121}z_{174}+z_{172}^{2}+z_{173}^{2}
\end{aligned}\]
\[\begin{aligned}Q_{145}={}&-z_{21}-z_{21}z_{176}-z_{21}z_{177}-z_{21}z_{183}-z_{21}z_{188}-z_{21}z_{193}
\\&-z_{21}z_{224}+z_{178}
\end{aligned}\]
\[\begin{aligned}Q_{146}={}&z_{4}z_{172}+z_{171}z_{172}+4z_{172}-z_{173}-z_{175}z_{178}-z_{176}
\end{aligned}\]
\[\begin{aligned}Q_{147}={}&-z_{179}z_{180}+z_{181}
\end{aligned}\]
\[\begin{aligned}Q_{148}={}&-1-2z_{120}z_{181}-2z_{121}z_{181}+z_{179}^{2}+z_{180}^{2}
\end{aligned}\]
\[\begin{aligned}Q_{149}={}&z_{4}z_{179}+z_{171}z_{179}-z_{178}z_{182}+4z_{179}-z_{180}-z_{183}
\end{aligned}\]
\[\begin{aligned}Q_{150}={}&-z_{184}z_{185}+z_{186}
\end{aligned}\]
\[\begin{aligned}Q_{151}={}&-1-2z_{120}z_{186}-2z_{121}z_{186}+z_{184}^{2}+z_{185}^{2}
\end{aligned}\]
\[\begin{aligned}Q_{152}={}&z_{4}z_{184}+z_{171}z_{184}-z_{178}z_{187}+4z_{184}-z_{185}-z_{188}
\end{aligned}\]
\[\begin{aligned}Q_{153}={}&-z_{189}z_{190}+z_{191}
\end{aligned}\]
\[\begin{aligned}Q_{154}={}&-1-2z_{120}z_{191}-2z_{121}z_{191}+z_{189}^{2}+z_{190}^{2}
\end{aligned}\]
\[\begin{aligned}Q_{155}={}&z_{4}z_{189}+z_{171}z_{189}-z_{178}z_{192}+4z_{189}-z_{190}-z_{193}
\end{aligned}\]
\[\begin{aligned}Q_{156}={}&-z_{194}z_{195}+z_{196}
\end{aligned}\]
\[\begin{aligned}Q_{157}={}&-1-2z_{120}z_{196}-2z_{121}z_{196}+z_{194}^{2}+z_{195}^{2}
\end{aligned}\]
\[\begin{aligned}Q_{158}={}&-z_{21}-z_{21}z_{198}-z_{21}z_{199}-z_{21}z_{205}+z_{200}
\end{aligned}\]
\[\begin{aligned}Q_{159}={}&z_{4}z_{194}+z_{171}z_{194}+4z_{194}-z_{195}-z_{197}z_{200}-z_{198}
\end{aligned}\]
\[\begin{aligned}Q_{160}={}&-z_{201}z_{202}+z_{203}
\end{aligned}\]
\[\begin{aligned}Q_{161}={}&-1-2z_{120}z_{203}-2z_{121}z_{203}+z_{201}^{2}+z_{202}^{2}
\end{aligned}\]
\[\begin{aligned}Q_{162}={}&z_{4}z_{201}+z_{171}z_{201}-z_{200}z_{204}+4z_{201}-z_{202}-z_{205}
\end{aligned}\]
\[\begin{aligned}Q_{163}={}&-z_{206}z_{207}+z_{208}
\end{aligned}\]
\[\begin{aligned}Q_{164}={}&-1-2z_{120}z_{208}-2z_{121}z_{208}+z_{206}^{2}+z_{207}^{2}
\end{aligned}\]
\[\begin{aligned}Q_{165}={}&z_{4}z_{206}-z_{126}z_{209}+z_{171}z_{206}+4z_{206}-z_{207}-z_{210}
\end{aligned}\]
\[\begin{aligned}Q_{166}={}&-z_{212}z_{213}+z_{214}
\end{aligned}\]
\[\begin{aligned}Q_{167}={}&-1-2z_{120}z_{214}-2z_{121}z_{214}+z_{212}^{2}+z_{213}^{2}
\end{aligned}\]
\[\begin{aligned}Q_{168}={}&z_{4}z_{212}-z_{126}z_{215}+z_{171}z_{212}+4z_{212}-z_{213}-z_{216}
\end{aligned}\]
\[\begin{aligned}Q_{169}={}&-z_{217}z_{218}+z_{219}
\end{aligned}\]
\[\begin{aligned}Q_{170}={}&-1-2z_{120}z_{219}-2z_{121}z_{219}+z_{217}^{2}+z_{218}^{2}
\end{aligned}\]
\[\begin{aligned}Q_{171}={}&-z_{21}-z_{21}z_{169}-z_{21}z_{170}-z_{21}z_{221}+z_{222}
\end{aligned}\]
\[\begin{aligned}Q_{172}={}&z_{4}z_{217}+z_{171}z_{217}+4z_{217}-z_{218}-z_{220}z_{222}-z_{221}
\end{aligned}\]
\[\begin{aligned}Q_{173}={}&z_{4}z_{128}+z_{127}+4z_{128}+z_{128}z_{171}-z_{135}-z_{178}z_{223}
\\&-z_{224}
\end{aligned}\]
\[\begin{aligned}Q_{174}={}&1995+216z_{0}+6z_{0}^{2}+8z_{0}z_{1}+16z_{0}z_{2}+2z_{0}z_{3}
\\&+8z_{0}z_{9}-2z_{0}z_{143}-4z_{0}z_{148}-2z_{0}z_{165}+154z_{1}+5z_{1}^{2}
\\&+16z_{1}z_{2}+8z_{1}z_{9}-4z_{1}z_{148}-2z_{1}z_{161}+308z_{2}+17z_{2}^{2}
\\&+18z_{2}z_{9}-8z_{2}z_{148}-2z_{2}z_{157}+32z_{3}+z_{3}^{2}-2z_{3}z_{165}
\\&+156z_{9}+5z_{9}^{2}-4z_{9}z_{148}-2z_{9}z_{157}+3z_{12}+z_{12}^{2}
\\&-2z_{12}z_{172}+15z_{13}+z_{13}^{2}-2z_{13}z_{179}+2z_{14}+z_{14}^{2}
\\&-2z_{14}z_{184}+2z_{15}+z_{15}^{2}-2z_{15}z_{189}+2z_{16}+z_{16}^{2}
\\&-2z_{16}z_{194}+2z_{17}+z_{17}^{2}-2z_{17}z_{201}+2z_{18}+z_{18}^{2}
\\&-2z_{18}z_{206}+2z_{19}+z_{19}^{2}-2z_{19}z_{212}+2z_{20}+z_{20}^{2}
\\&-2z_{20}z_{217}+4z_{23}+2z_{26}+z_{30}+z_{32}+z_{34}
\\&+z_{36}+2z_{40}+4z_{42}+2z_{43}+z_{46}+2z_{68}
\\&+z_{120}z_{225}+z_{121}z_{225}-z_{128}-32z_{143}+z_{143}^{2}-76z_{148}
\\&+z_{148}^{2}-4z_{157}+z_{157}^{2}-2z_{161}+z_{161}^{2}-32z_{165}
\\&+z_{165}^{2}-2z_{172}+z_{172}^{2}-2z_{179}+z_{179}^{2}-2z_{184}
\\&+z_{184}^{2}-2z_{189}+z_{189}^{2}-2z_{194}+z_{194}^{2}-2z_{201}
\\&+z_{201}^{2}-2z_{206}+z_{206}^{2}-2z_{212}+z_{212}^{2}-2z_{217}
\\&+z_{217}^{2}-z_{225}
\end{aligned}\]
\endgroup


\section{Complete compact circuit and fixed coefficients}\label{app:constants}

The register table below gives all 87 paid binary operations of the
lower-center equation. For the sharp-center equation replace the single
row \texttt{d=Am1*UE} by \texttt{d0=Am1*UE} followed by
\texttt{d=d0+UE}; every other row is unchanged. Its output
is the polynomial~\eqref{eq:compact}; the register names introduce no new
unknowns. The four ports are fixed by the finite recipe following the
table. The recipe is transcribed from the pinned complete parent display,
and its full four-integer output is independently bound by hashes and
bit lengths in the circuit artifact.
\begingroup\small
\begin{longtable}{@{}rll@{}}
\toprule Step & Register & Binary definition\\\midrule\endhead
1 & \texttt{bJ} & \texttt{beta} \texttt{*} \texttt{J}\\
2 & \texttt{q} & \texttt{bJ} \texttt{+} \texttt{1}\\
3 & \texttt{qmF} & \texttt{q} \texttt{-} \texttt{F}\\
4 & \texttt{C} & \texttt{qmF} \texttt{-} \texttt{alpha}\\
5 & \texttt{X} & \texttt{w} \texttt{*} \texttt{q}\\
6 & \texttt{k} & \texttt{eta} \texttt{+} \texttt{zeta}\\
7 & \texttt{kY} & \texttt{k} \texttt{*} \texttt{Y}\\
8 & \texttt{c} & \texttt{kY} \texttt{+} \texttt{eta}\\
9 & \texttt{q2} & \texttt{q} \texttt{*} \texttt{q}\\
10 & \texttt{mu} & \texttt{q2} \texttt{+} \texttt{1}\\
11 & \texttt{Xp1} & \texttt{X} \texttt{+} \texttt{1}\\
12 & \texttt{muY} & \texttt{mu} \texttt{*} \texttt{Y}\\
13 & \texttt{Am1} & \texttt{muY} \texttt{*} \texttt{Xp1}\\
14 & \texttt{4Am1} & \texttt{4} \texttt{*} \texttt{Am1}\\
15 & \texttt{Hp} & \texttt{4Am1} \texttt{-} \texttt{1}\\
16 & \texttt{Ap1} & \texttt{Am1} \texttt{+} \texttt{2}\\
17 & \texttt{muY2} & \texttt{muY} \texttt{*} \texttt{muY}\\
18 & \texttt{V0} & \texttt{X} \texttt{*} \texttt{muY2}\\
19 & \texttt{V0k} & \texttt{V0} \texttt{*} \texttt{k}\\
20 & \texttt{V0kk} & \texttt{V0k} \texttt{+} \texttt{k}\\
21 & \texttt{V0k2} & \texttt{V0k} \texttt{*} \texttt{V0kk}\\
22 & \texttt{Dtau} & \texttt{V0k2} \texttt{+} \texttt{1}\\
23 & \texttt{G1} & \texttt{bJ} \texttt{*} \texttt{qmF}\\
24 & \texttt{G} & \texttt{G1} \texttt{+} \texttt{alpha}\\
25 & \texttt{q2m1} & \texttt{q2} \texttt{-} \texttt{1}\\
26 & \texttt{Gq} & \texttt{G} \texttt{*} \texttt{q2m1}\\
27 & \texttt{qMF} & \texttt{q} \texttt{*} \texttt{MF}\\
28 & \texttt{MCqMF} & \texttt{MC} \texttt{+} \texttt{qMF}\\
29 & \texttt{RJ} & \texttt{MCqMF} \texttt{*} \texttt{J}\\
30 & \texttt{R} & \texttt{Gq} \texttt{+} \texttt{RJ}\\
31 & \texttt{q3} & \texttt{q2} \texttt{*} \texttt{q}\\
32 & \texttt{kR} & \texttt{k} \texttt{-} \texttt{R}\\
33 & \texttt{e} & \texttt{kR} \texttt{-} \texttt{1}\\
34 & \texttt{K0w} & \texttt{K0} \texttt{+} \texttt{w}\\
35 & \texttt{K0wC} & \texttt{K0w} \texttt{*} \texttt{C}\\
36 & \texttt{Bt} & \texttt{K0wC} \texttt{-} \texttt{F}\\
37 & \texttt{c2} & \texttt{c} \texttt{*} \texttt{c}\\
38 & \texttt{U} & \texttt{i} \texttt{*} \texttt{c2}\\
39 & \texttt{UE} & \texttt{U} \texttt{*} \texttt{Dtau}\\
40 & \texttt{Delta} & \texttt{Am1} \texttt{*} \texttt{Ap1}\\
41 & \texttt{SE} & \texttt{Delta} \texttt{*} \texttt{UE}\\
42 & \texttt{DpI2} & \texttt{SE} \texttt{*} \texttt{U}\\
43 & \texttt{Df} & \texttt{DpI2} \texttt{+} \texttt{1}\\
44 & \texttt{aF0} & \texttt{DpI2} \texttt{*} \texttt{Df}\\
45 & \texttt{fouraF0} & \texttt{4} \texttt{*} \texttt{aF0}\\
46 & \texttt{Fsqr} & \texttt{fouraF0} \texttt{+} \texttt{1}\\
47 & \texttt{Q} & \texttt{Delta} \texttt{*} \texttt{fouraF0}\\
48 & \texttt{Tp} & \texttt{T} \texttt{+} \texttt{alpha}\\
49 & \texttt{cT} & \texttt{c} \texttt{*} \texttt{Tp}\\
50 & \texttt{cTR} & \texttt{cT} \texttt{-} \texttt{R}\\
51 & \texttt{DfcTR} & \texttt{Fsqr} \texttt{*} \texttt{cTR}\\
52 & \texttt{V} & \texttt{DfcTR} \texttt{-} \texttt{c}\\
53 & \texttt{V2} & \texttt{V} \texttt{*} \texttt{V}\\
54 & \texttt{P1} & \texttt{q3} \texttt{*} \texttt{bJ}\\
55 & \texttt{mu2} & \texttt{mu} \texttt{*} \texttt{mu}\\
56 & \texttt{P1mu2} & \texttt{P1} \texttt{*} \texttt{mu2}\\
57 & \texttt{d} & \texttt{Am1} \texttt{*} \texttt{UE}\\
58 & \texttt{bB} & \texttt{cTR} \texttt{*} \texttt{Bt}\\
59 & \texttt{dmbB} & \texttt{d} \texttt{-} \texttt{bB}\\
60 & \texttt{Vcenter} & \texttt{V} \texttt{*} \texttt{dmbB}\\
61 & \texttt{center} & \texttt{n} \texttt{+} \texttt{Vcenter}\\
62 & \texttt{Hn} & \texttt{Hp} \texttt{*} \texttt{center}\\
63 & \texttt{threec} & \texttt{3} \texttt{*} \texttt{c}\\
64 & \texttt{twoX} & \texttt{2} \texttt{*} \texttt{X}\\
65 & \texttt{gap} & \texttt{threec} \texttt{-} \texttt{twoX}\\
66 & \texttt{Kpos0} & \texttt{X} \texttt{*} \texttt{gap}\\
67 & \texttt{Kpos} & \texttt{Kpos0} \texttt{+} \texttt{2}\\
68 & \texttt{lin} & \texttt{Hn} \texttt{+} \texttt{Kpos}\\
69 & \texttt{Zpart1} & \texttt{P1mu2} \texttt{*} \texttt{lin}\\
70 & \texttt{qm1Y} & \texttt{bJ} \texttt{*} \texttt{Y}\\
71 & \texttt{q3Bt} & \texttt{q3} \texttt{*} \texttt{Bt}\\
72 & \texttt{inner} & \texttt{qm1Y} \texttt{-} \texttt{q3Bt}\\
73 & \texttt{mu2inner} & \texttt{mu2} \texttt{*} \texttt{inner}\\
74 & \texttt{P1e} & \texttt{P1} \texttt{*} \texttt{e}\\
75 & \texttt{brk} & \texttt{mu2inner} \texttt{-} \texttt{P1e}\\
76 & \texttt{Zpart2} & \texttt{Hp} \texttt{*} \texttt{brk}\\
77 & \texttt{Znum} & \texttt{Zpart1} \texttt{+} \texttt{Zpart2}\\
78 & \texttt{Mtot} & \texttt{P1mu2} \texttt{*} \texttt{Hp}\\
79 & \texttt{M2} & \texttt{Mtot} \texttt{*} \texttt{Mtot}\\
80 & \texttt{Dg} & \texttt{Dtau} \texttt{*} \texttt{Df}\\
81 & \texttt{Z2} & \texttt{Znum} \texttt{*} \texttt{Znum}\\
82 & \texttt{RG} & \texttt{M2} \texttt{*} \texttt{Dg}\\
83 & \texttt{RGV2} & \texttt{RG} \texttt{*} \texttt{V2}\\
84 & \texttt{inside} & \texttt{Z2} \texttt{-} \texttt{RGV2}\\
85 & \texttt{Qinside} & \texttt{Q} \texttt{*} \texttt{inside}\\
86 & \texttt{correction} & \texttt{RG} \texttt{-} \texttt{Z2}\\
87 & \texttt{N} & \texttt{Qinside} \texttt{+} \texttt{correction}\\
\bottomrule\end{longtable}
\endgroup


The following finite Python recipe defines the four integers without a
repository checkout or a network request. Every range is fixed; no input
variable of the Diophantine equation occurs in it. The surrounding source
file also checks the resulting byte hashes against the pinned parent.
\lstinputlisting[language=Python]{report_inputs/constant_recipe.py}

\section{Artifact inventory and source references}

\begin{longtable}{@{}p{.35\textwidth}p{.60\textwidth}@{}}
\toprule Path & Purpose\\\midrule\endhead
\file{graph/equations/} & Four literal height-one polynomials, readable expansions, graph certificates and metrics.\\
\file{graph/independent_verify.py} & Independent finite algebra, graph and complete coefficient checker.\\
\file{graph/source324/}, \file{graph/source442/} & Exact pinned source equations; signed324 proof and calibration scope.\\
\file{direct/equations/} & Natural226 and natural227 residuals, complete expansions, projection certificates and integer678 recipe.\\
\file{direct/PROOF.md} & Detailed pooled-digit and fused-cap proof, including witness order.\\
\file{compact/} & Literal 87- and 88-operation polynomials, fixed coefficient reconstruction, parent comparison and exact checks.\\
\file{verification/} & Final integrated verification receipts and document review.\\
\file{PROVENANCE.json} & Inspected repository commit and component source boundaries.\\
\file{MANIFEST.sha256} & Integrity inventory of the delivered files.\\
\bottomrule
\end{longtable}

\begin{thebibliography}{9}
\bibitem{p127} ProveIt, package 127, \emph{Shared arithmetic certificates
and graph energies for Con(ZFC)}. Pinned
\href{https://github.com/VladimirReshetnikov/ProveIt/tree/1c5367452a0bd4dd7e9e4bceb1508b62dbd2ff61/Computability/DiophantineConsistency/con-zfc-diophantine-forms/127-graph-certificates}{repository package}.
The inherited graph energy, shared hub construction and sufficient
four-cycle condition are predecessors of the present compiler.
\bibitem{p125} ProveIt, package 125, \emph{Round 15 unified binomial and
even Pell index equations}. Pinned
\href{https://github.com/VladimirReshetnikov/ProveIt/tree/1c5367452a0bd4dd7e9e4bceb1508b62dbd2ff61/Computability/DiophantineConsistency/con-zfc-diophantine-forms/125-unified-binomial-round15}{repository package}.
Source of the signed442 comparison equation.
\bibitem{p128} ProveIt, package 128, \emph{One-row Pell certificates}.
Pinned \href{https://github.com/VladimirReshetnikov/ProveIt/tree/1c5367452a0bd4dd7e9e4bceb1508b62dbd2ff61/Computability/DiophantineConsistency/con-zfc-diophantine-forms/128-one-row-pell}{repository package}.
Predecessor for combining a cap and an index congruence.
\bibitem{p129} ProveIt, package 129, \emph{A semantic Pell anchor}.
Pinned \href{https://github.com/VladimirReshetnikov/ProveIt/tree/1c5367452a0bd4dd7e9e4bceb1508b62dbd2ff61/Computability/DiophantineConsistency/con-zfc-diophantine-forms/129-semantic-pell-anchor}{repository package}.
The source of the semantic label anchor, CRT lift, consecutive ordinates
and the natural233/natural234 comparison equations.
\bibitem{p131} ProveIt, package 131, \emph{Shared Pell caps, round 16}.
Pinned \href{https://github.com/VladimirReshetnikov/ProveIt/tree/1c5367452a0bd4dd7e9e4bceb1508b62dbd2ff61/Computability/DiophantineConsistency/con-zfc-diophantine-forms/131-shared-pell-caps-round16}{repository package}.
Source of the signed324 equation and its ordinary arithmetic calibration.
\bibitem{p133} ProveIt, package 133, \emph{Staged certificates}.
Pinned \href{https://github.com/VladimirReshetnikov/ProveIt/tree/1c5367452a0bd4dd7e9e4bceb1508b62dbd2ff61/Computability/DiophantineConsistency/con-zfc-diophantine-forms/133-staged-certificates}{repository package}.
The 84/436 and 85/328 parents, product-square proof and complete indexed
coefficient tuple of the compact component.
\bibitem{dprm} J.~Bayer, M.~David, B.~Stock, A.~Pal, Y.~Matiyasevich and
D.~Schleicher, \emph{Diophantine Equations and the DPRM Theorem}, Archive
of Formal Proofs. \url{https://isa-afp.org/entries/DPRM_Theorem.html}.
General background only; no unspecified DPRM polynomial is used as a
delivered equation in this report.
\end{thebibliography}
\end{document}
